\documentclass[11pt]{article}
\usepackage[T1]{fontenc}
\usepackage{lmodern}
\usepackage[margin=1in]{geometry}
\usepackage{amsmath,amssymb,amsthm,mathtools}
\usepackage{microtype,booktabs,longtable,array,enumitem}
\usepackage[dvipsnames]{xcolor}
\usepackage[obeyspaces,spaces]{xurl}
\usepackage{fancyhdr}
\usepackage[colorlinks=true,linkcolor=MidnightBlue,citecolor=MidnightBlue,urlcolor=MidnightBlue]{hyperref}
\usepackage{listings}
\lstset{basicstyle=\ttfamily\footnotesize,breaklines=true,columns=fullflexible,keepspaces=true,frame=single,rulecolor=\color{black!15}}
\newtheorem{theorem}{Theorem}[section]
\newtheorem{lemma}[theorem]{Lemma}
\newtheorem{proposition}[theorem]{Proposition}
\newtheorem{corollary}[theorem]{Corollary}
\theoremstyle{definition}
\newtheorem{definition}[theorem]{Definition}
\newtheorem{example}[theorem]{Example}
\newtheorem{remark}[theorem]{Remark}
\newcommand{\ipra}{\mathrm{iPRA}}
\newcommand{\PR}{\mathrm{PR}}
\newcommand{\N}{\mathbb N}
\newcommand{\Z}{\mathbb Z}
\newcommand{\Q}{\mathbb Q}
\newcommand{\C}{\mathbb C}
\newcommand{\bool}{\{0,1\}}
\newcommand{\dotminus}{\mathbin{\dot-}}
\newcommand{\len}{\operatorname{len}}
\newcommand{\Height}{\operatorname{ht}}
\newcommand{\Res}{\operatorname{Res}}
\newcommand{\lc}{\operatorname{lc}}
\newcommand{\Supp}{\operatorname{Supp}}
\newcommand{\Iter}{\operatorname{Iter}}
\newcommand{\Run}{\operatorname{Run}}
\newcommand{\Good}{\operatorname{Good}}
\newcommand{\coeff}{\operatorname{coeff}}
\newcommand{\code}[1]{\ifmmode\text{\nolinkurl{#1}}\else\nolinkurl{#1}\fi}
\DeclareMathOperator{\Lex}{Lex}
\setlength{\emergencystretch}{2em}
\hypersetup{pdftitle={Finite mathematical interfaces in iPRA},pdfauthor={Research report prepared for Vladimir Reshetnikov},pdfsubject={Demand compilation, controlled descent, and effective plane elimination}}
\setlist{itemsep=3pt,topsep=5pt}
\pagestyle{fancy}
\fancyhf{}
\fancyhead[L]{\small Finite mathematical interfaces in iPRA}
\fancyhead[R]{\small Research report}
\fancyfoot[C]{\thepage}
\setlength{\headheight}{14pt}
\title{\textbf{Finite mathematical interfaces in iPRA}\\[5pt]
\large Demand compilation, controlled descent,\\and effective plane elimination}
\author{Research report prepared for Vladimir Reshetnikov}
\date{15 November 2026}
\begin{document}
\maketitle
\begin{abstract}
We develop a reusable discipline for expressing mathematical constructions in
first-order intuitionistic primitive recursive arithmetic. The starting point
is ProveIt's existing finite-observation work, rather than a replacement for
its kernel-checked presentation of arithmetic. The new work addresses three
specific gaps. We construct primitive-recursive syntax for a recursion whose
precision parameter changes in the recursive call, and provide a checked
term-to-symbol compiler with an actual embedded arithmetic correctness proof.
We give explicit primitive-recursive execution bounds for a fixed number of
lexicographically descending counters, even when lower counters reset and no
runtime bound was supplied in advance. Finally, bounded coordinate changes
and Sylvester determinants give a finite-list B\'ezout theorem and a complete
rational-point algorithm for certified plane intersections. A finite bound
on auxiliary equations then converts pointwise auxiliary witnesses into a
complete point list without requiring a chosen global cover.

All main mathematical arguments are developed explicitly. The archive
separates kernel-checked Lean results, concrete unproved Lean targets, written
object-theory proof transformations, and finite executable diagnostics.
General Siegel, Faltings, effective Roth, and the remaining prime-distribution
problems are not claimed solved. The framework identifies the exact extra
arithmetic data their proposed routes still require.
\end{abstract}
\vspace{4pt}
\noindent\textbf{Source baseline.} ProveIt commit
\code{b0f6f535ffc16324683fef2c263260a88ae03755}.
\par\noindent\textbf{Reading guide.} Sections~\ref{sec:framework}--\ref{sec:descent}
give the general method; Sections~\ref{sec:elimination}--\ref{sec:applications}
give the mathematical applications and their limits;
Section~\ref{sec:lean} gives the formalization path and exact proof status.
\tableofcontents
\newpage
\section{The problem exposed by the repository}
\label{sec:scope}

The root arithmetic README makes a decisive distinction. The
\code{ConciseIPRA} library supplies kernel-checked equivalences between
presentations of intuitionistic primitive recursive arithmetic. The
\code{docs/proofs} directory contains 60 delivered research packets,
including nine framework packets, whose claims were filed as delivered;
the index does not certify them. In particular, the root README says that
none of those packets supplies a completed derivation of its main
application in \code{Prf PRAAx}~\cite{repo-root,repo-index}.
This report uses that distinction throughout. An excellent mathematical
outline, a total Lean function, a successful exact computation, and an
embedded arithmetic proof are four different forms of evidence.

\subsection{What can actually be improved}

The earlier frameworks already supply bounded Booleanization, finite
contracts, precision budgets, conditional inverses, finite nets,
metastability, and several kinds of coherent approximation.
The finite-demand packet treats nonmonotone backward precision schedules
and aggregate root supports. The coherent-methods packet supplies filtered
lifting, Hurwitz coefficients, and weighted determinant norms. The
finite-observation packet contains degree-uniform approximate
factorization, convexity-based continuity bounds, and finite-net
minimization~\cite{repo-demand,repo-coherent,repo-observation}.
Repeating these methods under a new name would not resolve a new difficulty.

The present extension concentrates on the following three gaps.

\begin{enumerate}
\item \textbf{A missing executable syntax bridge.}
The earlier \code{BackwardDemand.lean} defines a semantic
\code{compiledRecursion} and proves its recurrence equations, but explicitly
leaves the construction of its \code{PR} syntax and the object-theory
compiler unimplemented. We supply actual syntax for that program, checked
semantic correctness of the generated syntax, and a general checked
term-to-symbol correctness theorem inside \code{Prf PRAAx}. The remaining
internal proofs are identified individually.
\item \textbf{Termination arguments with resets.}
A single decreasing natural number is often too restrictive for an
algorithm that finishes one task, lowers a more significant counter, and
restarts a lower counter. A fixed finite lexicographic rank suffices. We
derive its execution bound from a primitive-recursive finite simulation,
so a separate numerical runtime bound need not be guessed. The distinction
between a fixed number of counter levels and an input-dependent number is
essential.
\item \textbf{General plane intersections and auxiliary equations.}
The earlier zero-count repairs use monic plane norms, hyperelliptic
curves, or particular weighted leading forms. A bounded shear removes the
coordinate restriction. A Sylvester determinant supplies both the sharp
degree-product count and a complete rational-point procedure for every
intersection passing its nonvanishing test. A finite coefficient envelope
then upgrades pointwise auxiliary witnesses into a uniform complete list.
\end{enumerate}

These contributions extend standard ideas in primitive recursion,
termination analysis, and elimination theory; no claim of historical
priority is made. The gain is an explicit constructive formulation,
quantifier and resource accounting, an executable syntax bridge, and a
precise comparison with the inspected repository proposals.

\subsection{A targeted comparison}

\begin{longtable}{>{\raggedright\arraybackslash}p{0.22\textwidth}>{\raggedright\arraybackslash}p{0.33\textwidth}>{\raggedright\arraybackslash}p{0.37\textwidth}}
\toprule
Existing packet & Reused achievement & Addition or retained boundary\\
\midrule\endhead
\code{ipra_embedding_framework}
& Finite observations; simple analytic branches; hyperelliptic zero counts
& Arbitrary plane intersections passing a resultant test; complete rational lists.\\
\code{ipra_embedding_framework_2026-11-15}
& Finite traces, guarded operations, finite nets, height/list distinction
& A resource hierarchy that constructs fuel for fixed-depth ranked runs.\\
\code{ipra_framework}
& Finite contracts; explicit power-gap families; object-theory infrastructure
& No repetition of the power-gap application; polynomial elimination supplies a different finite-list producer.\\
\code{ipra_framework_2026-11-15}
& Certificates, moduli, elementary bounded descent
& Resetting counters and fixed-versus-variable depth are made explicit.\\
\code{ipra_finite_demand}
& Backward demands; semantic parameter-changing recursion; root supports
& Actual \code{PR} syntax for its recursion. Exact singular Puiseux coefficients and multiplicity matching remain outside this report.\\
\code{ipra_finite_observation}
& Variable-depth tables; complete approximate factorization; restricted curve section bounds
& Finite auxiliary equations yield effective point lists when their coefficient heights are bounded.\\
\code{ipra_coherent_methods_2026-11-15}
& Filtered lifting, coherent towers, weighted norms
& No weight or smoothness assumption in the plane-intersection theorem.\\
\code{ipra_evidence_framework}
& Relational finite outputs and finite compatibility
& Finite exhaustive envelopes remove a pointwise auxiliary choice; no compactness is inferred.\\
\code{ipra_uniform_certificates}
& Uniform finite-net and root certificates
& A syntax-level bridge and a separate termination resource layer.\\
\bottomrule
\end{longtable}

The other four framework notes are indexed together in~\cite{repo-other};
the earlier embedding and contract packages are also consulted directly
\cite{repo-embedding,repo-framework}.
This is a targeted comparison of framework sources and selected difficulty
reports. It is not a proof audit of all 60 application packets.
The provenance inventory records exact paths and hashes for the inspected
sources. Where a prior result is needed, its status remains visible; a
README's self-description is not silently promoted to a verified theorem.

\subsection{The object theory and the ambient metatheory}

We use the standard first-order presentation of \(\ipra\): intuitionistic
first-order logic with equality; symbols for primitive-recursive
definitions, including composition and primitive recursion, with their
defining equations; successor injectivity and separation from zero; and
induction for quantifier-free formulas. Quantifier-free includes the
intuitionistic propositional connectives. Bounded formulas below are
admissible abbreviations after the Booleanization proof in
Section~\ref{sec:framework}.

The metatheory is ordinary mathematics, and Lean is used to check selected
metatheorems and to construct selected object proofs. Thus the following
statements must not be confused:
\[
  \text{Lean proves }\forall n\,P(n),
  \qquad
  \text{Lean constructs }\code{Prf PRAAx []}\,\ulcorner\forall n\,P(n)\urcorner.
\]
Only the second statement directly establishes derivability in the
embedded arithmetic theory. A semantic proof can help discover a finite
arithmetic proof, but no general reflection principle is available to
convert it automatically.

Similarly, \(\forall f\) in a theorem about all primitive-recursive
functions normally belongs to the metatheory: for each fixed definition
of \(f\), construct an arithmetic symbol and proof. It is not a
first-order quantification over an internal type of total functions.
Numerical parameters of a fixed definition, matrix sizes, polynomial
degrees, and lengths of finite lists may vary inside arithmetic.
The permitted variation must be established by the program construction,
not inferred merely from the notation.

\section{A general framework of finite interfaces and resources}
\label{sec:framework}

The method has two layers. A mathematical object is described by finite
observations and laws about them. A computation is justified by an
arithmetic program, a finite demand for the observations it uses, and a
bounded invariant. The program's termination resource is kept separate
from its mathematical correctness. This separation lets the same
representation support a predetermined recursion, an adaptive ranked
search, or an exhaustive certificate enumeration.

\subsection{Coding conventions}

Natural numbers are the only object-theory sort. Signed integers may be
coded by \(0\mapsto0\), \(n>0\mapsto2n\), and
\(-n\mapsto2n-1\). A rational is a pair consisting of an integer numerator
and a positive denominator; reduce by a bounded Euclidean algorithm when
a canonical code is needed. Invalid denominator codes are assigned a
fixed default value, while every mathematical theorem retains the
positive-denominator guard. Complex rational numbers are pairs of
rationals. The coding of finite lists supplies total primitive-recursive
length, lookup, append, and bounded fold functions. Values outside a
list's range have a default value and are never used without a range
test in a mathematical contract.

For definiteness, pairing can be coded by
\(\langle a,b\rangle=2^a(2b+1)-1\). Its projections use bounded search
below the code plus one. A list can be a length and a nested pairing,
with list traversal bounded by that length. Every concrete loop below
has a primitive-recursive bound; efficient binary encodings can replace
these without changing the arguments. Rational arithmetic, coefficient
arrays, determinant permutation sums, polynomial substitution, and
finite coefficient comparison are consequently primitive recursive.
There is no appeal here to a universal evaluator for arbitrary programs.

\begin{lemma}[Boolean compilation of fixed bounded formulas]
\label{lem:boolean}
For each fixed bounded formula \(A(\vec x)\) over primitive-recursive
symbols, there is a primitive-recursive \(a(\vec x)\in\bool\) such that
\(\ipra\) proves
\[
             A(\vec x)\ \longleftrightarrow\ a(\vec x)=1.
\]
The construction also supplies proofs of the range condition and of the
equivalence. It is a transformation on a fixed formula, not an internal
truth predicate for all arithmetic formulas.
\end{lemma}

\begin{proof}
Equality has a primitive-recursive comparison test, whose correctness
follows from the defining equations and induction on subtraction or
simultaneous predecessor comparisons. Atomic predicates are treated by
their fixed characteristic definitions. For Boolean values, conjunction
is multiplication, disjunction is the capped sum, and negation is
\(1-a\). Since the induction hypothesis proves the subformulas
decidable, implication is represented by
\(\max(1-a,b)\). This use of decidability is restricted to these
subformulas; no excluded-middle schema is introduced.

For \(\forall i<t(\vec x)\,B(i,\vec x)\), define the prefix product
\[
 u(0,\vec x)=1,\qquad
 u(n+1,\vec x)=u(n,\vec x)b(n,\vec x).
\]
Quantifier-free induction proves that the product remains Boolean.
To obtain a particular factor from \(u(n)=1\), fix its index \(i\)
as a parameter and induct on \(n\) with the quantifier-free formula
\(u(n)=1\to(i<n\to b(i)=1)\). In the successor case, a product
of two Boolean factors equals one only when both factors equal one;
compare the fixed index with the preceding prefix length.
For the converse, retain \(\forall i<t\,B(i)\) in the fixed context;
induct on \(n\) with \(n\le t\to u(n)=1\), using just the instance
at \(i=n\) in the successor case. This proves the required equivalence
without induction on the bounded universal formula itself.

For a bounded existential, use the capped prefix sum, or the
primitive-recursive first-success selector with a success bit. An
induction on the scan length maintains the Boolean equation saying
that the success bit agrees with the capped sum, and, when successful,
that the selected index is below the scan length and satisfies the
characteristic test. The forward existential implication instantiates
the product or sum at a supplied witness; the reverse implication uses
the selected index. The formula induction is a finite metatheoretic
construction, and all numerical inductions it emits are quantifier-free.
\end{proof}

\begin{corollary}[Bounded induction and bounded choice]
\label{cor:bounded}
Induction on any fixed bounded formula is derivable in \(\ipra\).
If \(\exists y\le B(\vec x)\,R(\vec x,y)\) is proved with bounded
\(R\), the bounded first-success scan supplies a fixed
primitive-recursive witness with an internal correctness proof.
Finite products of such scans give finite choice over bounded domains.
\end{corollary}

\begin{proof}
Replace the induction formula by the equation in
Lemma~\ref{lem:boolean}, perform quantifier-free induction, and replace
back. For choice, use the scan invariant just proved; the supplied
existence theorem forces its success bit to be one. Default outputs
make the function total even when the premise is false.
\end{proof}

\subsection{Finite observations}

\begin{definition}[An arithmetic observation interface]
Fix a primitive-recursive producer \(O(\vec a,k)\) of finite codes,
a bounded validity predicate \(V(\vec a,k,z)\), and bounded relations
for comparing or restricting observations. The parameters \(\vec a\)
describe an input presentation; \(k\) is an observation demand, not
necessarily a decimal or binary precision. An object assertion is an
ordinary first-order formula built from these contracts, for example
\[
   \forall k\,V(\vec a,k,O(\vec a,k)).
\]
No object of an infinite set or function space is thereby added to the
arithmetic domain.
\end{definition}

An observation may be a jet through degree \(k\), a matrix entry table,
a rational approximation with error \(2^{-k}\), a finite graph, a finite
list of roots without labels, or a polynomial identity certificate.
An equality of represented objects can be a family of bounded
observational equalities. An apartness assertion must carry an explicit
positive lower bound when an inverse is computed.

For instance, a real presentation can use a fixed rational-valued
program \(A(\vec a,k)\) and the regularity contract
\[
 |A(\vec a,p)-A(\vec a,q)|\le 2^{-p}+2^{-q}.
\]
Each instance is an integer inequality after clearing positive
denominators. One may reason under its universal closure, but cannot use
that universal closure as a quantifier-free induction formula. For a
fixed requested accuracy, the proof records exactly the finitely many
instances required. An equality test between arbitrary such names is
not a Boolean operation; approximate comparison with a specified margin
or a supplied exact algebraic code is a different interface.

\subsection{Operations and proof-bearing contracts}

An operation consists of fixed PR syntax for its numerical output,
fixed PR syntax for its finite input demands, and a derivation of its
local contract. For one input a typical contract is
\begin{equation}\label{eq:local-operation}
 V_A(\vec a,D(\vec a,k),z)
 \ \longrightarrow\
 V_B(\vec a,k,T(\vec a,k,z)).
\end{equation}
If several predecessor observations are needed, use a bounded list of
queries and the bounded conjunction of their contracts. If the producer
needs one demand for computing its output and a larger one for checking
the output, record both. Neither demand is inferred from the fact that
an ambient mathematical operation exists.

\begin{proposition}[Proof-producing composition]
\label{prop:composition}
Two operations of the form \eqref{eq:local-operation} compose. The
composite demand is obtained by substituting the second demand into the
first; the composite output is ordinary PR composition. Substitution
and implication elimination compose the two supplied internal proofs.
For a finite family of queries, a prefix maximum covers all demanded
indices, provided the validity interface explicitly supports the
corresponding prefix restriction.
\end{proposition}

\begin{proof}
If \(D_1\) supplies the input needed to produce an intermediate
observation at demand \(r\), and \(D_2\) is the demand for the final
request \(k\), use the first contract at \(r=D_2(k)\). Its conclusion
is exactly the premise of the second contract. PR syntax is closed
under the displayed substitutions. These are operations on given proof
terms; they do not use the semantic truth of an arbitrary input
certificate. For multiple queries the finite maximum dominates each
listed index by its defining recursion. Apply the restriction proof
before the local contract; absent that proof, a maximum is not a valid
replacement for several nonmonotone demands.
\end{proof}

\subsection{Execution resources and bounded invariants}

\begin{lemma}[Guarded finite execution]
\label{lem:execution}
Let \(u(0)=u_0\) and \(u(t+1)=T(\vec a,t,u(t))\) be an ordinary
primitive recursion, and let \(B(\vec a)\) be PR. Suppose a fixed context
\(\Gamma\), independent of the induction variable, proves a bounded
invariant \(I(\vec a,0,u_0)\) and the bounded-range step
\[
 t<B(\vec a),\ I(\vec a,t,u(t))
       \ \longrightarrow\ I(\vec a,t+1,u(t+1)).
\]
Then \(\Gamma\) proves \(I(\vec a,B(\vec a),u(B(\vec a)))\).
The context itself need not consist of bounded formulas.
\end{lemma}

\begin{proof}
Booleanize \(I\), keeping \(\vec a\) fixed. Induct on
\(t\le B(\vec a)\to I(\vec a,t,u(t))\), expressed by the corresponding
quantifier-free Boolean implication. From \(t+1\le B\) obtain \(t<B\)
and the preceding invariant; outside that range the implication has no
obligation. Instantiate the resulting universal formula at \(t=B\).
\end{proof}

There are three ways to supply \(B\) in this report. A bounded search
comes with its enumeration bound. A predetermined observation recursion
comes with its requested terminal stage, while its changing observation
parameters are compiled backwards. An adaptive computation with a
fixed-dimensional lexicographic rank receives a bound from
Section~\ref{sec:descent}. This third case permits output-dependent
control flow, provided the transition itself is fixed PR and the local
rank and invariant proofs are supplied.

\begin{theorem}[Restricted compilation principle]
\label{thm:compilation}
Consider a finite mathematical development whose representations and
operations have the above interfaces. Allow PR composition, bounded
folds and scans, the predetermined demand recursions of
Section~\ref{sec:demand}, and the fixed-depth ranked executions of
Section~\ref{sec:descent}. Require internal local proofs for each
primitive contract, proofs of initialization and preservation of a bounded
invariant for every numerical loop, and a proof that its terminal
invariant implies the requested output contract.
Then every output contract obtained by these composition, induction, and
output rules has an \(\ipra\) derivation under the supplied input
hypotheses. Every output witness
has fixed PR syntax. The transformation is a metatheoretic construction
from this finite development and its local derivations.
\end{theorem}

\begin{proof}
Proceed through the finite construction. Composition uses
Proposition~\ref{prop:composition}; scans and folds use
Corollary~\ref{cor:bounded} and Lemma~\ref{lem:execution}; a changing
parameter is replaced by the explicit two-recursion program of
Section~\ref{sec:demand}; a ranked loop is replaced by a fixed-fuel
simulation using Section~\ref{sec:descent}. Every generated invariant
is bounded, and Lemma~\ref{lem:boolean} supplies its quantifier-free
form. Existential introduction uses the explicitly generated witness
term. Universal introduction over numerical parameters is performed
only after their finite proof has been constructed and only when those
parameters do not occur freely in the retained context (otherwise the
relevant assumptions must first be discharged). No other induction
principle, soundness reflection, or uniform evaluator is used.
\end{proof}

The principle is deliberately restricted. It does not say that every
true specification of a PR program is provable in \(\ipra\), nor that
every Lean program is PR, nor that all constructive mathematics falls
inside arithmetic. Those claims would conceal exactly the obligations
the framework is meant to expose.

\subsection{A useful test for a proposed embedding}

The following four questions are a compact mathematical review tool.
What finite code is returned? Which observations of the input justify
it? What fixed PR construction bounds its execution? Which bounded
invariant proves its contract? A fifth question is needed for a
completeness theorem: why must every solution be represented by a code
the finite procedure examines? Sections~\ref{sec:elimination} and
\ref{sec:applications} answer the fifth question using elimination
polynomials and bounded auxiliary equations. Mere cardinality does not
answer it.

\section{Compiling a changing observation parameter into PR syntax}
\label{sec:demand}

The finite-demand packet already gives the mathematical construction in
this section~\cite{repo-demand}. Its explicit outstanding task was to
construct the \code{PR} syntax and connect it to object arithmetic.
We reproduce the argument so that this report is self-contained, then
describe the new checked implementation. The distinction between the
inherited theorem and the new formalization is intentional.

\subsection{The equations that need compilation}

Let \(D:\N^2\to\N\), \(b:\N\to\N\), and
\(F:\N^3\to\N\) be fixed PR functions. We want
\begin{align}
 X(0,k)&=b(k),\label{eq:demand-zero}\\
 X(n+1,k)&=F\bigl(n,k,X(n,D(n,k))\bigr).
 \label{eq:demand-successor}
\end{align}
The recursive call changes \(k\). It is therefore not the ordinary PR
constructor with \(k\) held fixed. Treating \(X(n,\cdot)\) as an
accumulator of function type would also leave the first-order PR algebra.
Nevertheless, the future requests are determined by \(D\) before the
recursive answers are known. This gives a direct first-order solution.

For terminal data \((N,K)\), define
\begin{align}
 B(0;N,K)&=K,\nonumber\\
 B(j+1;N,K)&=D\bigl(N\dotminus(j+1),B(j;N,K)\bigr),
 \label{eq:backward}\\
 H(t;N,K)&=B(N\dotminus t;N,K).
 \label{eq:horizon}
\end{align}
Truncated subtraction makes the definitions total at every input.
Only \(t\le N\) is used in the correctness argument.
Then perform a forward recursion with the terminal data fixed:
\begin{align}
 V(0;N,K)&=b(H(0;N,K)),\nonumber\\
 V(t+1;N,K)&=F\bigl(t,H(t+1;N,K),V(t;N,K)\bigr),
 \label{eq:forward}\\
 X(N,K)&=V(N;N,K).
 \label{eq:compiled}
\end{align}
Both passes use ordinary primitive recursion. The forward pass may
recompute a demanded horizon instead of storing a whole table; either
choice has fixed PR syntax.

\begin{lemma}[Horizon identities]
\label{lem:horizon}
The definitions satisfy
\[
 H(N;N,K)=K,\qquad
 H(t;N,K)=D(t,H(t+1;N,K))\quad(t<N).
\]
They also satisfy the coherence identities
\begin{align}
 B(j+1;N+1,K)&=B(j;N,D(N,K)),\label{eq:back-coherent}\\
 H(t;N+1,K)&=H(t;N,D(N,K))\quad(t\le N).
 \label{eq:horizon-coherent}
\end{align}
These equations have quantifier-free induction proofs in \(\ipra\).
\end{lemma}

\begin{proof}
The terminal equation follows from \(N\dotminus N=0\) and the base
equation for \(B\). If \(t<N\), then
\[
 N\dotminus t=(N\dotminus(t+1))+1,
 \quad
 N\dotminus((N\dotminus(t+1))+1)=t.
\]
Substitute these into the successor equation for \(B\).

For \eqref{eq:back-coherent}, induct on \(j\). At \(j=0\), the left
side is \(D(N,K)\), which is the right base value. In the successor
case use the previous equality and the elementary identity
\[
 (N+1)\dotminus(j+2)=N\dotminus(j+1).
\]
This identity is valid even beyond the requested horizon. The resulting
equations have identical outer \(D\) arguments. Finally, for \(t\le N\),
\((N+1)\dotminus t=(N\dotminus t)+1\); substituting this value for
\(j+1\) in \eqref{eq:back-coherent} proves
\eqref{eq:horizon-coherent}. All numerical side conditions are Boolean
comparisons, and the induction formulas are equations or guarded equations.
\end{proof}

\begin{theorem}[Predetermined demand recursion]
\label{thm:demand}
The fixed PR function \(X\) of \eqref{eq:compiled} satisfies
\eqref{eq:demand-zero} and \eqref{eq:demand-successor}. Its construction
and the two equations have \(\ipra\) proofs for each fixed triple of
input PR definitions.
\end{theorem}

\begin{proof}
At \(N=0\), the horizon is \(K\) and the forward base is \(b(K)\).
For the successor equation compare the forward executions for
\((N+1,K)\) and \((N,D(N,K))\). Their initial values are equal by
\eqref{eq:horizon-coherent}. At every step \(t<N\), their horizon
values at \(t+1\) are equal by the same equation. Quantifier-free
induction on the guarded numerical equality
\[
 t\le N\ \longrightarrow\
 V(t;N+1,K)=V(t;N,D(N,K))
\]
therefore makes their values equal through time \(N\). The final step
of the longer execution is
\begin{align*}
 X(N+1,K)
 &=F(N,H(N+1;N+1,K),V(N;N+1,K))\\
 &=F(N,K,V(N;N,D(N,K)))\\
 &=F(N,K,X(N,D(N,K))).
\end{align*}
The middle equality is the terminal horizon equation and the guarded
forward equality just proved. No induction hypothesis quantifies over
all demands or over an infinite family of completed objects.
\end{proof}

\subsection{Correctness at a requested observation}

Suppose \(C(n,k)\) is bounded and \(\ipra\) proves
\[
 C(0,H(0;N,K)),\qquad
 n<N,\ C(n,D(n,k))\ \longrightarrow\ C(n+1,k)
\]
under a fixed context. Booleanize \(C\), hold \(N,K\) fixed, and
induct on
\[
             t\le N\ \longrightarrow\ C(t,H(t;N,K)).
\]
Lemma~\ref{lem:horizon} makes the previous observation exactly the
premise required by the local step. At \(t=N\), the terminal horizon
is \(K\), giving \(C(N,K)\). If the base theorem is available for
every demand and the retained context has no free \(N,K\), universal introduction over the independently fixed
\(N,K\) yields \(\forall N\forall K\,C(N,K)\).
The induction itself never uses \(\forall k\,C(t,k)\) as its invariant.

For example, if \(D(n,k)=2k+1\), elementary induction gives
\[
            H(t;N,K)=2^{N-t}(K+1)-1\qquad(t\le N).
\]
The resulting exponential precision demand is harmless for PR
definability. If \(D\) is nonmonotone, the original backward recurrence
still works; no maximum or monotonicity assertion is substituted for
its exact values. Several simultaneous queries are handled by the
bounded-prefix construction from the earlier packet, provided the
prefix validity interface is used explicitly.

\subsection{The new term-to-symbol compiler}

The repository represents a primitive-recursive function by an inductive
type \code{PR k}, with constructors for zero, successor, projections,
composition, and primitive recursion. A term \code{PTm k} may contain
variables and applications of these symbols. Define a compiler
\(S:\code{PTm k}\to\code{PR k}\) by
\[
 S(x_i)=\operatorname{proj}_i,
 \qquad
 S(f(t_0,\ldots,t_{m-1}))
   =\operatorname{comp}(f,S(t_0),\ldots,S(t_{m-1})).
\]

\begin{theorem}[Internal correctness of term compilation]
\label{thm:term-compile}
For every term \(t\), arbitrary tuple of substituted terms \(\sigma\),
and arbitrary context \(\Gamma\), there is an embedded arithmetic proof
\[
 \Gamma\vdash_{\mathrm{PRAAx}}
      \operatorname{app}(S(t),\sigma)=t[\sigma].
\]
The supplied Lean theorem \code{TermCompiler.symbol_correct} constructs
this proof and was checked by the Lean kernel.
\end{theorem}

\begin{proof}
Induct on the finite syntax of \(t\) in the metatheory. For a variable,
the projection axiom gives the equation. For an application, the
composition axiom gives
\[
 \operatorname{app}(S(t),\sigma)
   =f(\operatorname{app}(S(t_0),\sigma),\ldots,
       \operatorname{app}(S(t_{m-1}),\sigma)).
\]
The induction hypotheses are already object proofs for its argument
equations. Equality congruence and transitivity replace those arguments
by \(t_i[\sigma]\), yielding the desired equation. This proof consumes
only defining axioms and equality rules; no theorem about the truth of
arbitrary denotations is used.
\end{proof}

\subsection{The concrete compiler and the exact completed boundary}

In \code{DemandCompiler.lean}, the input types are
\code{D : PR 2}, \code{b : PR 1}, and \code{F : PR 3}; the output is
\code{compiled D b F : PR 2}. The intermediate symbols are:

\begin{center}
\begin{tabular}{lll}
\toprule
Symbol & Arity & Arguments\\
\midrule
\code{backwardSymbol D} & 3 & reverse count, terminal time, demand\\
\code{horizonSymbol D} & 3 & terminal time, demand, forward time\\
\code{forwardSymbol D b F} & 3 & forward count, terminal time, demand\\
\code{compiled D b F} & 2 & terminal time, demand\\
\bottomrule
\end{tabular}
\end{center}

The primitive-recursion convention is the actual upstream convention:
the recursion counter comes first, and a step symbol receives
\((\text{index},\text{accumulator},\text{parameters})\).
Predecessor and truncated subtraction are themselves explicitly built
from \code{PR.prec}. This avoids any unverified claim that a host
function is primitive recursive.

The checked theorem \code{denote_compiled} identifies the generated
syntax with the earlier semantic two-pass program. The checked host
theorems \code{compiled_zero} and \code{compiled_succ} establish the
desired recurrence equations for its denotation. In addition,
\code{object_backward_zero} and \code{object_backward_succ} construct
actual \code{Prf PRAAx} proofs of the generic backward pass equations
for arbitrary fixed \(D\). They use Theorem~\ref{thm:term-compile} to
handle the composite step body.

The full object-theory versions of
\eqref{eq:demand-zero}--\eqref{eq:demand-successor} are stated as
concrete target propositions in \code{ObjectTargets.lean}; they are
\emph{not yet implemented as Lean proofs}. Their human proof is
Theorem~\ref{thm:demand}. The remaining Lean work is finite and explicit:
derive the required subtraction equations internally, prove backward
coherence by equation induction, prove guarded horizon and forward
coherence, and unfold the final composition. Section~\ref{sec:lean}
records the exact order and the upstream APIs involved.

\subsection{The boundary with nested adaptive demands}

The theorem permits a changing parameter because \(D(n,k)\) does not
depend on the as-yet-uncomputed value \(X(n,\cdot)\). It does not
license an arbitrary nested recursion in which the next recursive
argument is obtained from a previous recursive answer. For a branching
request tree whose child parameters are fixed PR functions of the
parent and whose depth is \(N\), a supplied PR branching and parameter
envelope makes bounded bottom-up evaluation available. For a genuinely
adaptive algorithm, a fixed PR transition and the lexicographic
criterion in the next section provide another route. Without such a
resource argument, a total host definition or a familiar termination
proof does not establish PR implementability.

\section{Primitive-recursive fuel for fixed-dimensional lexicographic descent}
\label{sec:lex-fuel}
\label{sec:descent}

This section addresses a recurring gap between an ordinary termination
argument and an argument usable in intuitionistic primitive recursive
arithmetic.  A process may have a decreasing pair or a decreasing finite
tuple of natural numbers, while a lower coordinate is reset to a much larger
value whenever a higher coordinate decreases.  Its termination is familiar,
but neither an induction on an existential termination assertion nor an
unbounded while-loop is automatically available in iPRA.  We give an
explicit primitive-recursive fuel function.  The mathematical ingredients
are standard fixed-depth nested recursion (compare the broader length-function analysis in~\cite{schmitz2014}); the useful point is the fully
specified interface and the automatic construction of its growth envelope.

Throughout this section, the dimension $d$ is a fixed external natural
number.  It is part of a theorem instance, not an unrestricted input to the
resulting primitive-recursive program.  Tuple operations, finite sequence
operations, bounded maxima, bounded searches, and finite iteration are
primitive recursive.  Their bounded assertions are represented by their
primitive-recursive Boolean characteristic functions, so the inductions
below have quantifier-free representatives in iPRA.

\subsection{A bound on controlled descending chains}

Write $v>_{\mathrm{lex}}w$ when, at the first coordinate at which the two
$d$-tuples differ, the coordinate of $v$ is larger.  The first coordinate
is the most significant one.  There are no strict comparisons between
empty tuples.

Let $h(x,t)$ be a primitive-recursive function, where $x$ denotes an
arbitrary finite tuple of numerical parameters.  Suppose that iPRA proves
that $h(x,t)$ is nondecreasing in $t$.  Define, separately for each fixed
$d$, the following primitive-recursive functions:
\begin{align}
 F_0(x,s)&=s+1,\\
 F_{d+1}(x,s)&=F_d(x,\cdot)^{[h(x,s)+1]}(s).
 \label{eq:lex-fuel-recursion}
\end{align}
Here $f^{[r]}(s)$ means $r$ successive applications of $f$ to $s$, with
$f^{[0]}(s)=s$.  Formula~\eqref{eq:lex-fuel-recursion} describes a sequence
of primitive-recursive definitions.  It does not assert that the
three-variable function $(d,x,s)\mapsto F_d(x,s)$ is primitive recursive.

\begin{lemma}[Monotonicity of the bounds]
For every fixed $d$, iPRA proves
\[
 F_d(x,s)\ge s+1,
 \qquad
 s\le t\ \Longrightarrow\ F_d(x,s)\le F_d(x,t).
\]
\end{lemma}

\begin{proof}
The assertion for $d=0$ is immediate.  Suppose it has been proved for
$d$.  Iterating a nondecreasing function preserves inequalities between
starting points.  Iterating a function $f(u)\ge u+1$ more times cannot
decrease its value, and one iteration already gives a value at least one
larger than its starting point.  These two assertions follow by
quantifier-free induction on the iteration counter.  If $s\le t$, then
$h(x,s)+1\le h(x,t)+1$, and hence
\[
 F_d(x,\cdot)^{[h(x,s)+1]}(s)
 \le F_d(x,\cdot)^{[h(x,s)+1]}(t)
 \le F_d(x,\cdot)^{[h(x,t)+1]}(t).
\]
There is at least one iteration in the definition of $F_{d+1}$, so its
successor bound follows as well.  The induction on $d$ here constructs a
separate finite iPRA proof for each fixed dimension.
\end{proof}

\begin{theorem}[Controlled lexicographic chain bound]
\label{thm:controlled-lex}
Fix $d$.  Let
\[
 v_s,v_{s+1},\ldots,v_{s+L-1}\in\mathbb N^d
\]
be a finite sequence such that $v_t>_{\mathrm{lex}}v_{t+1}$ whenever both
terms occur, and every coordinate of $v_t$ is at most $h(x,t)$.  Then
\[
                         s+L\le F_d(x,s).
\]
For each fixed $d$, this statement, with the finite sequence represented
by a numerical code, is provable in iPRA.
\end{theorem}

\begin{proof}
An empty sequence causes no difficulty.  If $d=0$, a strictly descending
sequence of empty tuples contains at most one term, and $s+L\le s+1$.

For the passage from $d$ to $d+1$, partition the sequence into its
maximal consecutive blocks on which the first coordinate is constant.
The first coordinate never increases, and it strictly decreases at
every boundary between blocks.  Its initial value is at most $h(x,s)$.
Consequently there are at most $h(x,s)+1$ blocks.

Let the block boundaries, as exclusive endpoints, be
\[
 s=s_0<s_1<\cdots<s_p=s+L.
\]
Within the $i$th block the remaining $d$ coordinates form a strictly
lexicographically descending sequence, with the same coordinate envelope.
The already established $d$-dimensional theorem therefore gives
\[
                         s_{i+1}\le F_d(x,s_i).
\]
By monotonicity and induction on the number of blocks,
\[
 s_i\le F_d(x,\cdot)^{[i]}(s).
\]
Since $p\le h(x,s)+1$, this gives
\[
 s+L=s_p
 \le F_d(x,\cdot)^{[p]}(s)
 \le F_d(x,\cdot)^{[h(x,s)+1]}(s)
 =F_{d+1}(x,s).
\]

For the iPRA implementation, a block boundary is the first index at
which the first coordinate changes, found by bounded search within the
given sequence.  A block subcode is obtained by a primitive-recursive
slice operation.  The block-count inequality is proved by induction on
the bounded list of boundaries: after $i$ boundaries, the initial first
coordinate is at least the current first coordinate plus $i$.  All
conditions on sequence positions and block numbers are bounded and are
compiled into quantifier-free predicates.  No induction on an assertion
of the form ``there exists a terminating computation'' is used.
\end{proof}

The bound is exact for the stated class of chains.  This is useful for
understanding why the dimension restriction matters.  For $d=0$, use the
one-term sequence of empty tuples.  Having constructed a maximal
$d$-dimensional chain starting at any time $t$, make a $(d+1)$-dimensional
chain with $h(x,s)+1$ blocks.  In block $i$, use first coordinate
$h(x,s)-i$ and the maximal $d$-dimensional suffix chain starting at
$F_d(x,\cdot)^{[i]}(s)$.  The envelope holds because $h$ is nondecreasing;
the first coordinate drops at each block boundary; and the final
exclusive endpoint is exactly $F_{d+1}(x,s)$.  For each fixed dimension
this construction is itself primitive recursive.

\subsection{From a transition system to a bounded witness}

Let $\tau(x,u)$ be a total primitive-recursive one-step transition on
numerically coded states, and let $u_0(x)$ be the initial state.  Define
the finite simulation before proving termination:
\[
 u(x,0)=u_0(x),\qquad
 u(x,t+1)=\tau(x,u(x,t)).
\]
This is ordinary primitive recursion.  A finite simulation remains
defined even when the corresponding unbounded while-loop would fail to
halt.  Let $A(x,u)$ be a primitive-recursive activity predicate, and make
inactive states absorbing.  Let $I(x,u)$ be a quantifier-free invariant
whose initialization and preservation have iPRA proofs.  Let
\[
 r(x,u)=(r_0(x,u),\ldots,r_{d-1}(x,u))
\]
be a fixed tuple of primitive-recursive rank functions.  Assume iPRA
proves the local obligation
\begin{align*}
 I(x,u)\wedge A(x,u)\wedge A(x,\tau(x,u))
 \ \Longrightarrow\ 
 r(x,u)>_{\mathrm{lex}}r(x,\tau(x,u)).
\end{align*}

\begin{theorem}[Automatic fuel for fixed-dimensional descent]
\label{thm:automatic-lex-fuel}
Under these assumptions there is an explicitly defined
primitive-recursive $B(x)$ such that iPRA proves
\[
                         \neg A(x,u(x,B(x))).
\]
If iPRA additionally proves
$I(x,u)\wedge\neg A(x,u)\to R(x,\operatorname{out}(x,u))$,
then $w(x)=\operatorname{out}(x,u(x,B(x)))$ is a
primitive-recursive witness and iPRA proves $R(x,w(x))$.
\end{theorem}

\begin{proof}
Define the trajectory envelope
\[
 h(x,t)=\max\bigl(\{0\}\cup
 \{r_j(x,u(x,q)):q\le t,\ j<d\}\bigr).
\]
This is a bounded maximum of explicitly given primitive-recursive
values, so it is primitive recursive.  Its monotonicity and the
coordinate-envelope property are proved by quantifier-free induction on
the bounded maximum.  Set $B(x)=F_d(x,0)$ using this $h$.

Quantifier-free induction on $t$ establishes $I(x,u(x,t))$.  Suppose
that $u(x,B(x))$ were active.  Since inactive states are absorbing, all
states $u(x,0),\ldots,u(x,B(x))$ would then be active.  This finite
backward implication may be proved by contraposition of the direct
forward persistence of inactivity; activity is decidable, so this uses
only decidable reasoning.  The local rank obligation gives a controlled
strictly descending chain of $B(x)+1$ rank tuples.  Theorem
\ref{thm:controlled-lex} would imply
$B(x)+1\le F_d(x,0)=B(x)$, a contradiction.  Therefore the final state
is inactive.  The postcondition follows by applying the given local
output theorem to this state and its invariant.

There is no circularity in the envelope: $u(x,t)$ executes an already
total one-step map for an externally supplied finite number of steps.
Its definition does not inspect a termination proof.  The fixed-depth
definitions of $F_d$ then compose and iterate this particular
primitive-recursive $h$.  The resulting bound can be extremely large;
its purpose is a primitive-recursive proof of totality, not an efficient
runtime estimate.
\end{proof}

This theorem is a proof transformer, not an internal reflection principle.
For a particular transition system it consumes iPRA derivations of the
invariant and local rank conditions and produces a derivation of its
bounded termination and postcondition.  Merely proving these conditions
in Lean's ambient logic does not by itself establish their derivability
in iPRA.

\subsection{A reset example}

Consider a state $(a,b,z)\in\mathbb N^3$ and the following transition:
\[
 \tau(a,b,z)=
 \begin{cases}
 (a,b,z),&a=b=0,\\
 (a,b-1,z+1),&b>0,\\
 (a-1,2^z,z+1),&b=0<a.
 \end{cases}
\]
The active predicate is $a+b>0$.  The pair $(a,b)$ strictly decreases
lexicographically at every active step, although $a+b$ can increase
drastically at a reset.  The unused coordinate $z$ records elapsed work
and can make later resets much larger than earlier ones.

With initial data $(a_0,b_0,z_0)$, direct quantifier-free induction gives
\[
 a_t\le a_0,\qquad z_t\le z_0+t,\qquad
 b_t\le b_0+2^{z_0+t}.
\]
Thus one can use the explicit monotone envelope
$h(t)=a_0+b_0+2^{z_0+t}$ instead of the automatic maximum.  The theorem
provides primitive-recursive fuel without an induction on a termination
formula.

For this particular example one can also calculate the exact runtime.
First consume the initial $b_0$ countdown steps.  Put
\[
 T_0=b_0,\qquad T_{i+1}=T_i+1+2^{z_0+T_i}.
\]
At time $T_i$, for $i\le a_0$, the state is
$(a_0-i,0,z_0+T_i)$.  If $i<a_0$, one reset step and then exactly
$2^{z_0+T_i}$ countdown steps establish the next displayed state.
Consequently the process halts at $T_{a_0}$.  These assertions are proved
by quantifier-free inductions on the countdown length and then on $i$.
For $b_0=z_0=0$, the first runtimes as $a_0$ varies are
$0,2,7,136,2^{136}+137$.  The rapid growth is still primitive recursive:
$i\mapsto T_i$ is a fixed primitive recursion.

There cannot be a natural-valued rank that decreases under every
lexicographic decrease in $\mathbb N^2$.  Such a rank $\rho$ would satisfy
$\rho(0,b)\ge b$ and
$\rho(1,0)>\rho(0,b)$ for every $b$, which is impossible at
$b=\rho(1,0)$.  The theorem therefore does not flatten the whole order
$\omega^2$ into $\mathbb N$.  It uses the particular transition system,
its input, and its controlled finite trajectory to obtain a fuel bound.

\subsection{Why dimension must remain fixed}

The restriction is substantive, even for a simple coordinate envelope.
Take $h(t)=t+1$, suppress $x$, and write
\[
 G_0(t)=t+1,\qquad G_{d+1}(t)=G_d^{[t+2]}(t).
\]
Each fixed $G_d$ is primitive recursive, nondecreasing, and at least the
successor function.  Also $G_{d+1}(t)>G_d(t)$.

\begin{proposition}[The uniform diagonal is not primitive recursive]
The function $D(n)=G_n(n)$ is not primitive recursive.
\end{proposition}

\begin{proof}
We first prove by structural induction on primitive-recursive
definitions that, for every primitive-recursive function $f$ of any
fixed finite arity, there are fixed numbers $m,N$ such that
\[
 t\ge N,\quad x_1,\ldots,x_k\le t
 \quad\Longrightarrow\quad f(x_1,\ldots,x_k)\le G_m(t).
 \tag{$*$}
\]
Zero, successor, and projections satisfy this property.  Increase the
indices and thresholds to a common value whenever finitely many
functions are involved; this is legitimate because $G_m(t)$ is
nondecreasing in $m$ as well as in $t$.

For composition, if the component functions and the outer function are
majorized by $G_m$, the composite is majorized by
$G_m(G_m(t))\le G_{m+1}(t)$.

For primitive recursion, suppose
$f(0,\vec x)=g(\vec x)$ and
$f(n+1,\vec x)=q(n,f(n,\vec x),\vec x)$, with $g$ and $q$ majorized by
$G_m$ above a common threshold.  Fix $t$ above that threshold and
$\vec x\le t$.  Induction on $i\le t$ gives
\[
 f(i,\vec x)\le G_m^{[i+1]}(t).
\]
Indeed, the arguments of $q$ at the next step are bounded by the current
iterate, which is at least $t$.  For $n\le t$ this yields
\[
 f(n,\vec x)\le G_m^{[t+1]}(t)\le G_{m+1}(t).
\]
This proves $(*)$ for all primitive-recursive definitions.

If $D$ were primitive recursive, $(*)$ would provide fixed $m,N$ with
$D(n)\le G_m(n)$ for all sufficiently large $n$.  Choose
$n\ge\max(N,m+1)$.  Monotonicity in the level and strict growth between
successive levels give
\[
 D(n)=G_n(n)\ge G_{m+1}(n)>G_m(n),
\]
a contradiction.
\end{proof}

Consequently a proof that every fixed-dimensional instance has
primitive-recursive fuel does not furnish a primitive-recursive fuel
function whose inputs include arbitrary dimension.  In particular, an
implementation that represents the rank by a list of input-dependent
length needs an additional argument.  This distinction remains even if
Lean accepts a total definition of $(d,t)\mapsto G_d(t)$ by recursion
whose intermediate values are functions.  Lean's acceptance of such a
definition is not a certificate that it belongs to the first-order
primitive-recursive function algebra.

In fact, there is no primitive-recursive function of $(d,s)$ that bounds
the lengths of \emph{all} chains with this envelope and arbitrary
dimension.  Exactness would give
$G_d(s)\le s+B(d,s)$ for any such bound $B$.  The primitive-recursive
diagonal $n\mapsto n+B(n,n)$ would then majorize $G_n(n)$, contrary to
the same majorization argument.  Particular families with a
variable-length rank can still admit better bounds because of
additional structure; the conclusion concerns the unrestricted
lexicographic criterion.

\subsection{Lean implementation obligations}

The semantic development can use fixed tuples, a Boolean lexicographic
comparison, finite lists, and the iterator.  Its essential lemmas are:
\begin{enumerate}
 \item monotonicity and the successor lower bound for each $F_d$;
 \item the bounded block decomposition of a lexicographic chain;
 \item the controlled-chain bound, by external induction on $d$;
 \item primitive recursiveness of finite simulation, running maxima, and
       $F_d$ for each fixed $d$;
 \item preservation of a quantifier-free invariant under finite
       simulation;
 \item the contradiction from $F_d(x,0)+1$ active rank states;
 \item the output postcondition for the resulting bounded execution.
\end{enumerate}
The important typing distinction is between the defensible statement
\[
 \texttt{\textup{for every }d,\ Primrec\ (fun\ s\ =>\ fuel\ d\ s)}
\]
and the false general assertion
\[
 \texttt{Primrec\ (fun\ (d,s)\ =>\ fuel\ d\ s)}.
\]
For the object-theoretic layer, encode each fixed tuple and finite list
by natural numbers, compile each bounded assertion to a Boolean
predicate, and attach iPRA derivations of the local invariant and rank
conditions.  The fixed-dimensional construction then yields object
derivations, rather than an inference from semantic truth in Lean to
provability in iPRA.

\section{Elimination certificates that produce complete affine rational-point lists}
\label{sec:elimination}

The preceding proof-transport discipline becomes particularly useful when
a classical finiteness argument has supplied auxiliary equations.  Existing
packets already contain complete rational-root scans
\cite{repo-geometry}, determinant norms for monic plane equations
\cite{repo-demand}, and sharp weighted refinements for several curve
families \cite{repo-coherent}.  We combine those finite mechanisms in a
different way: an arbitrary plane presentation is put into a suitable
coordinate system by a bounded search, and its elimination polynomial
produces a \emph{complete list} of rational intersection points.  A uniform
coefficient envelope on auxiliary equations then suffices to produce a
complete list for the target set, without choosing an auxiliary equation
for each point.  This makes explicit a stronger interface than the
finite-tuple covering premise in the Siegel packet \cite{repo-siegel}.

The algebra is standard elimination theory in origin
\cite{sturmfels}.  The contribution here is the full arithmetic
construction, its exact guards, its quantifier-sensitive composition, and
its induction audit.  No general Siegel or Faltings theorem is asserted.

\subsection{Exact inputs, guards, and bounded coordinates}

A bivariate rational polynomial is a finite coefficient array
\(F=(f_{ij})\), with default value zero outside the encoded array.  Its
value is the finite sum
\[
 F(x,y)=\sum_{i,j}f_{ij}x^iy^j.
\]
Rationals have canonical signed numerators and positive denominators.
Equality, order, polynomial evaluation, and every array operation used
below are primitive recursive.  The total degree of a nonzero polynomial
is the largest \(i+j\) having \(f_{ij}\ne0\); a separate Boolean flag
recognizes the zero polynomial.  All degrees below are total degrees
unless the variable is displayed explicitly.

The final algorithm is total and returns either a successful list or an
\emph{unsupported} flag.  If either polynomial is a nonzero constant, it
returns the empty list successfully.  If one is the zero polynomial and
neither is a nonzero constant, it returns unsupported.  Thus its main
branch has nonzero polynomials \(F,G\) of degrees \(d,e\ge1\).
The unsupported flag is never interpreted as an empty solution set.

For an integer \(c\ge0\), make the shear
\[
 U=X+cY,\qquad V=Y,\qquad
 F_c(U,V)=F(U-cV,V),\quad G_c(U,V)=G(U-cV,V).
 \tag{E1}\label{eq:elim-shear}
\]
The inverse map is \((U,V)\mapsto(U-cV,V)\); both maps preserve rational
coordinates and are mutually inverse by direct substitution.  Binomial
expansion computes the coefficient arrays.

\begin{lemma}[A bounded generic direction]\label{lem:elim-shear}
Let \(F_d,G_e\) be the nonzero highest homogeneous parts of \(F,G\).
There is a \(c\in\{0,\ldots,d+e\}\) such that
\[
 a(c):=F_d(-c,1)\ne0,\qquad b(c):=G_e(-c,1)\ne0.
 \tag{E2}\label{eq:elim-leading}
\]
The first such \(c\) is a primitive recursive function of the coefficient
arrays.  For that choice,
\[
 F_c=\sum_{i=0}^d f_i(U)V^i,\quad
 G_c=\sum_{j=0}^e g_j(U)V^j,
\]
where
\[
 f_d=a(c),\quad g_e=b(c),\quad
 \deg_U f_i\le d-i,\quad \deg_U g_j\le e-j
 \tag{E3}\label{eq:elim-degree-rows}
\]
for nonzero coefficients.  In particular, every rational specialization
\(F_c(u,V)\) has degree exactly \(d\).
\end{lemma}

\begin{proof}
The coefficients of the polynomial \(a(T)=F_d(-T,1)\) are the coefficients
of \(F_d\), with their signs changed according to the exponent of \(X\).
Thus \(a\) is nonzero and has degree at most \(d\).  Similarly \(b\) is
nonzero of degree at most \(e\).  Their product is nonzero: the product of
their leading coefficients is nonzero in \(\Q\).

We recall the finite root bound in the form needed here.  If a nonzero
degree-\(r\) rational polynomial vanishes at distinct supplied numbers
\(t_1,\ldots,t_s\), synthetic division gives
\(P(T)=(T-t_1)P_1(T)\).  Evaluating at \(t_j\ne t_1\) shows
\(P_1(t_j)=0\), because \(t_j-t_1\ne0\).  Continue along the finite
list.  A nonzero polynomial cannot lose more than its degree in these
divisions, so \(s\le r\).  Each quotient and division identity is finite
coefficient data; no search for an unspecified root is involved.

Consequently \(a(T)b(T)\) cannot vanish at all \(d+e+1\) candidates
\(0,\ldots,d+e\).  A scan with a success flag, frozen at the first passing
candidate, returns the required \(c\).  Finally the shear preserves the
bound on total degree.  A term containing \(V^d\) in \(F_c\) cannot
contain a positive power of \(U\), and its coefficient is
\(F_d(-c,1)\).  Every coefficient of \(V^i\) has \(U\)-degree at most
\(d-i\).  This proves all the assertions.
\end{proof}

Henceforth \(c\) always denotes this canonical first passing value.  The
constant leading coefficient is useful twice: it controls the elimination
degree, and it prevents the polynomial in the second coordinate from
vanishing identically at an exceptional value of the first coordinate.

\subsection{The finite elimination certificate and its degree}

Put \(N=d+e\).  Let \(S(U)\) be the \(N\)-by-\(N\) matrix whose rows
are coefficient vectors, in columns \(V^0,\ldots,V^{N-1}\), of
\[
 F_c,VF_c,\ldots,V^{e-1}F_c,
 \qquad G_c,VG_c,\ldots,V^{d-1}G_c.
\]
Define the elimination polynomial
\[
 R(U)=\det S(U).
 \tag{E4}\label{eq:elim-resultant}
\]
This is a Sylvester determinant; other conventional row orders change at
most its sign.  We use the permutation-sum definition, so the certificate
can be checked by finite rational polynomial arithmetic.  In the main
branch, \(\Good(F,G)\) means that the coefficient array of \(R\) is
nonzero.  It is a decidable primitive recursive test.  It is not presumed
to succeed on every pair.

\begin{lemma}[The total-degree audit]\label{lem:elim-degree}
If \(R\ne0\), then \(\deg R\le de\).
\end{lemma}

\begin{proof}
Consider one nonzero term of the determinant expansion.  In the row
\(V^rF_c\), an entry selected from column \(k\) is \(f_{k-r}\), and
its degree is at most \(d-(k-r)\).  In the row \(V^sG_c\), the
corresponding bound is \(e-(k-s)\).  A permutation uses each column
\(0,\ldots,d+e-1\) once.  Therefore the degree of the selected product
is at most
\begin{align*}
 2de-\sum_{k=0}^{d+e-1}k
       +\sum_{r=0}^{e-1}r+\sum_{s=0}^{d-1}s
 &=2de-\frac{(d+e)(d+e-1)}2
          +\frac{e(e-1)+d(d-1)}2\\
 &=de.
\end{align*}
A zero selected entry makes the determinant term zero and requires no
degree calculation.  Finite addition cannot create a coefficient above
the maximum support bound.  This proves the claim without any assumption
that leading determinant terms do not cancel.
\end{proof}

\begin{lemma}[A common point gives an elimination root]
\label{lem:elim-root}
If \(F(x,y)=G(x,y)=0\), then \(R(x+cy)=0\).
\end{lemma}

\begin{proof}
Write \(u=x+cy\), \(v=y\).  The column
\(w=(1,v,\ldots,v^{N-1})^{\mathsf T}\) satisfies \(S(u)w=0\): each
row evaluates a shifted copy of \(F_c\) or \(G_c\) at \(v\).
Multiply by the adjugate of \(S(u)\).  The finite adjugate identity gives
\(R(u)w=0\), and its first component is \(R(u)=0\).
Determinant specialization and the adjugate identity follow from their
finite polynomial expansions, so no assertion about a completed
algebraic closure is used.
\end{proof}

For complete enumeration this is already enough.  For the sharp
cardinality bound one must account for several points with the same
\(u\)-coordinate.  Counting only distinct roots of \(R\) would omit
that information.

\begin{lemma}[A fibre contributes its full multiplicity]
\label{lem:elim-fibre}
If \(v_1,\ldots,v_r\) are distinct rationals and
\(F_c(u,v_j)=G_c(u,v_j)=0\) for every \(j\), then
\[
 (U-u)^r\mid R(U).
 \tag{E5}\label{eq:elim-fibre-power}
\]
\end{lemma}

\begin{proof}
The polynomial \(F_c(u,V)\) has degree exactly \(d\), so the finite
root bound gives \(r\le d\le N\).  Form a constant \(N\)-by-\(N\)
matrix \(W\).  Its first \(r\) columns are
\[
 (1,v_j,\ldots,v_j^{N-1})^{\mathsf T},\qquad 1\le j\le r;
\]
its remaining columns are the standard basis columns
\(e_r,\ldots,e_{N-1}\), numbering coordinates from zero.  The
upper-left block is a Vandermonde matrix, the upper-right block is zero,
and the lower-right block is the identity.  Thus
\[
 \det W=\prod_{i<j}(v_j-v_i)\ne0.
\]
The Vandermonde formula is itself a finite polynomial identity: as a
polynomial in the last ordinate its determinant has the previous
ordinates as roots and the preceding Vandermonde determinant as its
leading coefficient; induction proves the displayed product.

Each of the first \(r\) columns of \(S(U)W\) vanishes entrywise at
\(U=u\).  Synthetic division therefore factors \(U-u\) from every
entry in each such column.  Factoring a scalar polynomial from these
\(r\) columns in the permutation-sum determinant proves
\((U-u)^r\mid\det(S(U)W)\).  Multiplicativity of the finite determinant
gives \(\det(S(U)W)=R(U)\det W\).  Dividing the quotient polynomial by
the nonzero rational \(\det W\) proves \eqref{eq:elim-fibre-power}.
\end{proof}

\begin{theorem}[Finite arithmetic B\'ezout bound]\label{thm:elim-bezout}
Assume \(\Good(F,G)\).  Every finite list of pairwise distinct rational
common zeros of \(F,G\) has length at most \(de\).
\end{theorem}

\begin{proof}
Apply the shear to the list and group entries by their first coordinate,
using decidable equality of canonical rationals.  Let the distinct first
coordinates be \(u_1,\ldots,u_t\), with group sizes
\(r_1,\ldots,r_t\).  Within a group the second coordinates are distinct,
because the shear is injective.  Lemma~\ref{lem:elim-fibre} gives
\((U-u_j)^{r_j}\mid R\) for every \(j\).

For completeness, coprimeness of these powers has explicit finite
witnesses.  If \(a\ne b\), put
\[
 L=\frac{U-a}{b-a},\qquad M=-\frac{U-b}{b-a},\qquad L+M=1.
\]
Expand \(1=(L+M)^{r+s-1}\).  A term whose exponent of \(L\) is at
least \(r\) is divisible by \((U-a)^r\); every other term has exponent
of \(M\) at least \(s\) and is divisible by \((U-b)^s\).  Grouping
the terms gives polynomial B\'ezout coefficients for the two powers.
Multiplying such identities supplies coprime cancellation for a finite
product.  Hence
\[
 \prod_{j=1}^t(U-u_j)^{r_j}\mid R(U).
\]
The product is monic of degree \(\sum_jr_j\).  Since \(R\ne0\),
degree comparison and Lemma~\ref{lem:elim-degree} yield
\[
 \operatorname{length}(\text{list})
     =\sum_jr_j\le\deg R\le de.
\]
All factorization witnesses and grouped lists are finite coded data.
\end{proof}

The statement needs neither an irreducibility oracle nor an unproved
claim about common components.  In classical algebra, absence of a common
nonconstant factor implies the good-shear resultant is nonzero.  That
characterization is unnecessary here: \(R\ne0\) is checked directly.
When it fails, even a finite or empty rational intersection may be
unsupported.  Completeness is asserted on the explicitly recognized
successful domain, not for all pairs of equations.

\subsection{From elimination to a complete list and height bound}

We include the rational-root argument to make the resulting positive
finiteness theorem self-contained.  This ingredient is already available
in the repository \cite{repo-geometry}; the new step is its composition
with the uniformly chosen two-coordinate elimination.

\begin{lemma}[Complete bounded rational-root scan]
\label{lem:elim-rational-roots}
For each nonzero rational polynomial \(P(T)\), a fixed primitive
recursive program returns exactly its rational roots, without duplicates.
\end{lemma}

\begin{proof}
Multiply all coefficients by their positive denominator product.  This
gives a nonzero integer polynomial with the same roots.  By bounded scans
find its least and greatest nonzero coefficient indices \(k,n\).  If
\(k>0\), record the root zero.  Write
\[
 P(T)=T^kQ(T),\qquad
 Q(T)=a_0+a_1T+\cdots+a_mT^m,
 \quad a_0a_m\ne0,quad m=n-k.
\]
If \(m=0\), there are no further roots.  Otherwise take a nonzero root
\(p/q\) in lowest terms, with \(q>0\).  Multiplication by \(q^m\)
gives the integer identity
\[
 a_mp^m+a_{m-1}p^{m-1}q+\cdots+a_0q^m=0.
 \tag{E6}\label{eq:elim-root-identity}
\]
It follows that \(p\mid a_0q^m\) and \(q\mid a_mp^m\).
The Euclidean algorithm gives a B\'ezout equation for the coprime
\(p,q\).  If an integer is coprime to two integers separately, multiplying
their B\'ezout equations proves it coprime to their product; induction
therefore proves coprimeness with the required powers.  Coprime
cancellation in \eqref{eq:elim-root-identity} now gives
\[
 p\mid a_0,\qquad q\mid a_m,
 \qquad 0<|p|\le|a_0|,quad 1\le q\le|a_m|.
 \tag{E7}\label{eq:elim-root-box}
\]
Scan this finite integer rectangle, retain coprime pairs whose exact
substitution into \(P\) is zero, and canonicalize.  The final equation
test proves soundness, and \eqref{eq:elim-root-box} proves that every
nonzero rational root is tested.  The separately recorded zero is a root
exactly when \(k>0\).  Canonical deduplication preserves membership.

A nonzero constant returns the empty list.  The zero polynomial has no
finite root-list contract and receives a failure flag in the total
wrapper.  Every scan bound is a numerical function of already computed
coefficients, so all operations are primitive recursive.
\end{proof}

\begin{theorem}[A positive plane-intersection theorem]
\label{thm:elim-complete}
There is a fixed primitive recursive function \(\operatorname{Intersect}\)
such that, if \(\Good(F,G)\), its output is a duplicate-free list \(L\)
with
\[
 (x,y)\in L\quad\Longleftrightarrow\quad F(x,y)=G(x,y)=0
 \quad\text{for all }x,y\in\Q.
 \tag{E8}\label{eq:elim-complete}
\]
Moreover \(|L|\le de\).  It also computes an upper bound on the canonical
codes, and on the naive coordinate heights, of all these points.
\end{theorem}

\begin{proof}
The algorithm is explicit:
\begin{enumerate}
\item Compute the canonical shear \(c\) and its nonzero polynomial \(R\).
\item Enumerate all rational roots \(u\) of \(R\), using
Lemma~\ref{lem:elim-rational-roots}.
\item For each \(u\), enumerate all rational roots \(v\) of
\(F_c(u,V)\).  Its leading coefficient is the fixed nonzero rational
\(a(c)\), so this call never encounters the zero polynomial.
\item Retain the pairs satisfying \(G_c(u,v)=0\), replace them by
\((u-cv,v)\), and remove duplicates.
\end{enumerate}
Every retained point satisfies both original equations by the inverse
shear.  Conversely let \(F(x,y)=G(x,y)=0\).  Put \(u=x+cy\) and
\(v=y\).  Lemma~\ref{lem:elim-root} puts \(u\) in the first complete
root list.  The equation \(F_c(u,v)=0\) puts \(v\) in that fibre's
complete list.  The other equation passes the final filter, and the
inverse shear returns \((x,y)\).  This proves
\eqref{eq:elim-complete}.  Theorem~\ref{thm:elim-bezout} applies to the
duplicate-free output and gives its length bound.

For a canonical rational \(z=p/q\), put \(h(z)=\max(|p|,q)\), and
for a pair put \(h(x,y)=\max(h(x),h(y))\).  The maximum of these heights
over \(L\), with zero for an empty list, bounds every rational common
point by completeness.  Taking the maximum of the point codes works in
the same way.  Neither maximum is an unbounded search.

The outer root rectangle is determined by \(R\); each inner rectangle
is determined by the coefficients of \(F_c(u,V)\) already computed for
that outer entry.  Nesting loops with such primitive recursive bounds
again gives a primitive recursive program.  Thus the conclusion includes
an actual list and height producer, rather than only a cardinality bound.
\end{proof}

The theorem is a zero-dimensional result: successful elimination has
proved that the relevant intersection has a finite rational list.  It
does not bound all rational points on a curve in the absence of auxiliary
equations.

\subsection{Two cover interfaces and a bounded-envelope theorem}

Fix a nonzero nonconstant \(F\), and a fixed primitive recursive
decidable predicate \(A(x,y)\) on rational codes.  Ordinary numerical
parameters of this predicate may be carried through every program and
proof below; they are suppressed in the notation.  Integrality of both
coordinates is one example.  Define the target set by
\[
 S(x,y)\quad:\Longleftrightarrow\quad F(x,y)=0\ \wedge\ A(x,y).
\]
This is notation for a numerical predicate, not an additional set sort.

\begin{corollary}[A supplied global auxiliary cover]
\label{cor:elim-global-cover}
Suppose a finite list \(G_1,\ldots,G_m\) is supplied, each pair
\((F,G_i)\) passes the good-resultant test, and the following mathematical
premise has been proved:
\[
 \forall x,y\in\Q\quad
 S(x,y)\ \longrightarrow\ \bigvee_{i=1}^mG_i(x,y)=0.
 \tag{E9}\label{eq:elim-global-cover}
\]
Then a primitive recursive function of these data produces a complete
list of the target points and a height bound.  Its length is at most
\(\sum_i d\deg G_i\).
\end{corollary}

\begin{proof}
Concatenate the lists \(\operatorname{Intersect}(F,G_i)\), filter by
\(A\), and deduplicate.  A retained point satisfies \(F=0\) and \(A\).
If \(S(x,y)\), eliminate the supplied finite disjunction
\eqref{eq:elim-global-cover}.  In each branch the corresponding
intersection theorem puts \((x,y)\) in the union.  Its length is bounded
by the sum of the intersection length bounds; its height maximum is a
bound for all target points.

No finite computation is said to decide the universal cover premise.
That premise is used with a proof, or retained as a hypothesis.  The
program only computes the lists justified by the displayed contracts.
\end{proof}

A pointwise auxiliary-existence assertion suffices when all its witnesses
belong to one explicitly bounded coefficient space.  This observation is
useful because no choice function or globally selected cover is needed.

For \(e,H\ge1\), let \(\mathcal B(e,H)\) be the finite list of nonzero
rational polynomials of total degree at most \(e\), with every reduced
coefficient numerator bounded in magnitude by \(H\) and every positive
denominator bounded by \(H\).  Zero coefficients are represented by
\(0/1\).  There are
\[
 K=\frac{(e+1)(e+2)}2
\]
monomials of degree at most \(e\).  Listing all numerator-denominator
pairs and then all coefficient vectors gives at most
\[
 B(e,H)=\bigl((2H+1)H\bigr)^K
 \tag{E10}\label{eq:elim-box-size}
\]
raw vectors; normalization and removal of zero polynomials can only
decrease that number.  Nonzero constants may also be removed, since they
cannot vanish at a target point.  A radix-vector scan implements this
enumeration by one fixed primitive recursive definition even when \(e\)
is an input.

\begin{theorem}[A bounded auxiliary envelope gives positive finiteness]
\label{thm:elim-envelope}
Assume
\[
 \forall x,y\in\Q\quad S(x,y)\ \longrightarrow\
 \exists G\in\mathcal B(e,H)\,
       \bigl(\Good(F,G)\ \wedge\ G(x,y)=0\bigr).
 \tag{E11}\label{eq:elim-pointwise-envelope}
\]
There is a fixed primitive recursive algorithm, depending on the fixed
predicate \(A\), which takes \(F,e,H\) and produces a complete list of
target points and a height or code bound.  A coarse bound on the number
of points is
\[
 d e\, B(e,H).
 \tag{E12}\label{eq:elim-envelope-card}
\]
In particular, the conclusion does not require a selector taking a point
to its auxiliary polynomial.
\end{theorem}

\begin{proof}
Enumerate the entire finite coefficient box.  Retain every \(G\) for
which the good-resultant test succeeds, and concatenate all corresponding
lists \(\operatorname{Intersect}(F,G)\).  Filter the concatenation by
\(A\) and remove duplicates.  Every retained point satisfies \(S\).

Conversely, given \(S(x,y)\), eliminate the bounded existential in
\eqref{eq:elim-pointwise-envelope}.  Its witness occurs in the enumerated
box, passes the test, and has \(G(x,y)=0\).  The complete intersection
list for that witness therefore contains \((x,y)\).  Hence the union
contains it, and the true predicate \(A(x,y)\) retains it.  This proves
completeness without selecting witnesses as a function of points.

There are at most \(B(e,H)\) candidates and each successful candidate
has degree at most \(e\), so its list has at most \(de\) entries.
This proves \eqref{eq:elim-envelope-card}.  Taking a maximum over the
computed complete list gives the height or code bound.  The bound on the
coefficient box is explicit, and every subsequent list-producing call is
primitive recursive; their composition is primitive recursive as well.
\end{proof}

The degree \emph{and coefficient} bounds are doing different jobs.  The
degree controls intersection cardinality.  The coefficient bound makes
the space of all possible witnesses a finite searchable list, which
produces a height bound and complete enumeration.

For comparison, suppose only that every finite list of target points can
be covered by at most \(M\) auxiliary equations of degree at most \(e\),
each with a nonzero elimination certificate.  The auxiliaries may depend
on the whole finite list and their coefficients may be unbounded.
Assign every list entry to its first covering equation.  Each group has
at most \(de\) points by Theorem~\ref{thm:elim-bezout}, hence every
distinct target list has length at most \(Mde\).  This is a valid finite
cardinality argument.  It supplies neither a global cover to use in
Corollary~\ref{cor:elim-global-cover} nor a coefficient envelope to use
in Theorem~\ref{thm:elim-envelope}.  In particular, changing auxiliary
equations along larger and larger lists does not give a stopping rule for
enumerating all target points.  The missing uniformity has been located,
not inferred away.

\subsection{Exact examples and what the construction repairs}

Take
\[
 F=X^2+Y^2-5,\qquad G=XY-2.
\]
The first passing shear has \(c=1\), giving
\[
 F_c=U^2-2UV+2V^2-5,\qquad G_c=UV-V^2-2,
 \qquad R(U)=(U^2-9)^2.
\]
For \(u=3\) the retained ordinates are \(1,2\); for \(u=-3\) they
are \(-1,-2\).  The inverse shear returns precisely
\[
 (1,2),\ (2,1),\ (-1,-2),\ (-2,-1).
\]
There are four points, attaining \(de=4\), although only two first
coordinates after the canonical shear.  Thus fibre multiplicity is
necessary for the sharp count in this example.

For \(F=XY-2\) and \(G=X+Y-3\), the candidate \(c=0\) fails the
leading-coefficient test for \(F\), and \(c=1\) fails it for \(G\).
The scan chooses \(c=2\).  Here
\[
 F_c=UV-2V^2-2,\qquad G_c=U-V-3,
 \qquad R(U)=-(U-4)(U-5)
\]
with our row convention.  The complete list is \((1,2),(2,1)\),
attaining \(de=2\).  This illustrates a presentation in which the
bounded shear is genuinely needed.

If \(F=Y-X\) and \(G=2Y-2X\), the resultant is zero and the algorithm
returns unsupported.  Returning an empty list would be false: the
intersection contains every rational point on that line.  The code also
handles elimination roots at zero by stripping the initial polynomial
power before the rational-root scan.

Relative to the audited packets, the construction adds three concrete
interfaces: a bounded conversion from arbitrary plane coordinates to
constant leading coefficients, a complete rational-point producer for
every successful intersection certificate, and a finite-envelope route
from bounded auxiliary existence to positive finiteness.  It does not
provide the quantitative approximation theorem or coefficient bounds
needed to apply that route to all curves in the general Siegel or
Faltings statements.  It also does not provide actual rational
Mordell--Weil quotient representatives or effective exceptional-zero
estimates for prime-distribution arguments.

\subsection{Quantifier-free induction audit and Lean completion}

All claims above have a direct first-order arithmetic form over natural
codes.  A displayed bounded conjunction means the corresponding
primitive recursive Boolean fold equals one.  Validity guards on rational
and polynomial codes prevent division by zero; invalid input branches
receive fixed defaults.  The following are the main proof obligations.

\begin{enumerate}
\item \textbf{Polynomial operations.}  Implement coefficient arrays,
finite evaluation, multiplication, shear, and synthetic division.  Prove
the division identity by induction on the remaining coefficient count.
The invariant is a coefficientwise equality of coded polynomials, with
finite support and degree bounds compiled into Boolean tests.
\item \textbf{Determinants.}  Enumerate length-\(N\) vectors over
\(\{0,\ldots,N-1\}\), retain vectors with distinct entries, and sum
their signed products.  The loop bound \(N^N\) is primitive recursive.
Prove specialization, the adjugate identity, determinant multiplication,
and column factor extraction by finite expansion and reindexing.  The
degree audit records the sum of selected column indices; it needs no
unbounded algebraic assertion.
\item \textbf{Root scans.}  Compute denominator products, first and last
nonzero indices, gcds, and the rectangles
\eqref{eq:elim-root-box}.  The scan invariant says that every stored
entry passes the exact equation test and that every passing candidate in
the processed prefix is stored.  Both assertions are bounded Boolean
predicates.  Instantiate the separately proved divisibility bound after
the scan to obtain completeness for an arbitrary rational root.
\item \textbf{Intersection lists.}  The nested-loop invariant expresses
membership in the accumulated union of the processed finite root
fibres.  Completeness for an arbitrary point follows by two instances of
the root-list theorem and Lemma~\ref{lem:elim-root}.  The sharp length
bound additionally needs Vandermonde and the finite grouped-factor
invariant \(R=P_jQ_j\), with all quotient coefficients stored.
\item \textbf{Coefficient envelopes.}  The coefficient-box enumeration
has an explicit radix bound.  The union invariant identifies stored
points with passing intersections in the processed prefix.  The
unbounded cover premise is held in the proof context and instantiated
only after this induction.  It is never an induction formula.
\end{enumerate}

In Lean, the native specification may use rational coefficient lists and
finite arrays.  The essential definitions are
\texttt{shear}, \texttt{canonicalShear}, \texttt{sylvester},
\texttt{elimination}, \texttt{good}, \texttt{rationalRoots},
\texttt{intersect}, \texttt{auxiliaryBox}, and
\texttt{envelopeList}.  Define
\[
\begin{split}
 \operatorname{Common}(F,G,p)&:
       F(p_1,p_2)=0\wedge G(p_1,p_2)=0,\\
 \operatorname{Complete}(L,P)&:
       \forall p\,(p\in L\leftrightarrow P(p)).
\end{split}
\]
The principal native theorem targets are precisely the statements of
Lemmas~\ref{lem:elim-shear}--\ref{lem:elim-rational-roots} and
Theorems~\ref{thm:elim-complete}, \ref{thm:elim-envelope}.  The
complete-list theorem can be finished before the fibre-multiplicity
module: it uses the common-point root lemma, not the sharp cardinality
bound.  Thus that longer determinant argument need not block the positive
enumeration result.

The object-language phase uses the same fixed arithmetic programs,
their compiled Boolean predicates, and the displayed quantifier-free
invariants to construct actual \texttt{Prf PRAAx} derivations.  A native
Lean theorem about rational arrays is not yet such a derivation, and
\texttt{pra\_iff\_native} is applied only after that object proof has
been built.  The accompanying exact Python code checks the concrete
examples and exceptional branches; it is a reference implementation, not
a proof of the universal statements or an iPRA proof checker.

\section{Which difficulties are resolved, and which data are still missing}
\label{sec:applications}

The new results make two kinds of progress. They turn certain informal
resource arguments into actual PR programs, and they turn particular
geometric certificates into complete point lists. They do not supply
every certificate that an existing application would need. This section
states the resulting interfaces precisely.

\subsection{Finite exhaustive witnesses as a general interface}

The auxiliary-polynomial construction is an instance of a useful
general rule. It applies to finite chart families, bounded algebraic
certificates, bounded collections of exact matrices, and finite
geometric decompositions, provided each candidate has a certified
complete finite output.

\begin{theorem}[Elimination of a bounded auxiliary choice]
\label{thm:exhaustive}
Let \(P(a,p)\), \(\Good(a,c)\), and \(R(a,c,p)\) be fixed PR
predicates on natural codes. Let \(B(a)\) and the finite list
\(L(a,c)\) be fixed PR programs. Suppose \(\ipra\) proves, for every
good candidate,
\[
 \Good(a,c)\ \longrightarrow\
 \bigl(p\in L(a,c)\ \longleftrightarrow\ R(a,c,p)\bigr).
\]
Suppose also that it proves the pointwise coverage assertion
\begin{equation}\label{eq:bounded-witness-cover}
 P(a,p)\ \longrightarrow\
 \exists c\le B(a)\,
       \bigl(\Good(a,c)\wedge R(a,c,p)\bigr).
\end{equation}
Then a fixed PR program produces a duplicate-free complete list of
all \(p\) satisfying \(P(a,p)\), and \(\ipra\) proves its completeness.
No chosen function sending \(p\) to a candidate \(c\) is required.
\end{theorem}

\begin{proof}
Scan \(c=0,\ldots,B(a)\). For good candidates concatenate \(L(a,c)\),
retain precisely entries satisfying \(P(a,\cdot)\), and remove
duplicates by a bounded scan through the accumulated list.
All loops have PR bounds: the outer counter is explicit and each
inner list has a PR length. The scan invariant says, for a fixed
parameter \(p\), that membership in the accumulated list is
equivalent to
\[
 P(a,p)\wedge\exists c<i\,
       (\Good(a,c)\wedge p\in L(a,c)).
\]
This is a bounded formula; its quantifier-free compilation is the
induction formula. At the end, every retained entry satisfies \(P\)
by the filter. Conversely, given \(P(a,p)\), eliminate the existential
in \eqref{eq:bounded-witness-cover}. Its witness is an enumerated good
candidate, and the candidate-completeness theorem puts \(p\) into
the corresponding list. The invariant places it in the final list.
Universal introduction over \(p\) follows after the finite induction.
\end{proof}

If \(P\) is not decidable, omit the filter. The same argument yields
a finite list containing every target point, hence a height or code
bound, but not necessarily an exact list of target points. This is
often the correct output contract when the remaining property is
not an arithmetic Boolean test.

Theorem~\ref{thm:exhaustive} does not require an explicit bound on
the \emph{number} of successful candidates smaller than \(B(a)+1\).
It relies instead on a finite bound on their possible codes. This is
why a coefficient bound on auxiliary equations can do more than a
cardinality bound on exceptional subspaces.

\subsection{A precise obstruction to converting cardinality into a list}

The distinction above is not only a bookkeeping concern. There is an
elementary computability counterexample, already anticipated by the
repository's height/list warnings~\cite{repo-faltings,repo-framework}.

\begin{proposition}[Uniform singleton bounds do not locate their member]
\label{prop:singleton}
There is one fixed PR predicate \(P(e,t)\) such that \(\ipra\) proves
that every duplicate-free list satisfying \(P(e,\cdot)\) has length
at most one, uniformly in \(e\), but there is no total computable
procedure producing a complete finite list for every \(e\). Nor is
there a total computable bound \(H(e)\) such that
\(P(e,t)\to t\le H(e)\) for all \(e,t\).
\end{proposition}

\begin{proof}
Fix an ordinary effective enumeration of machines and let \(P(e,t)\)
say that machine \(e\) first halts at step \(t\). A bounded machine
simulation is PR, and ``first'' is a bounded test. If \(s<t\) and
\(P(e,s)\), then the negative earlier-halting condition in \(P(e,t)\)
fails. Decidable comparison of \(s,t\) therefore proves
\(P(e,s)\wedge P(e,t)\to s=t\) in \(\ipra\). Applying that fact
to two distinct positions proves the list bound.

A total computable complete-list procedure would decide whether the
machine halts by testing whether its output list is empty. A total
computable bound \(H(e)\) would decide the same question by running
the machine through that many steps and checking the finite trace.
Either conclusion contradicts the elementary undecidability of the
halting problem. The simulator is bounded; no total universal
primitive-recursive evaluator is being postulated.
\end{proof}

The example can be placed directly inside the auxiliary-polynomial
interface. Take the fixed curve \(F(X,Y)=Y\) and the target points
\((t,0)\) with \(P(e,t)\). Every target point admits the linear
auxiliary equation \(G_t(X,Y)=X-t\), and \(F,G_t\) have a nonzero
resultant. Every finite target list is covered by one such auxiliary
equation (use \(X\) for an empty list). Degrees are uniformly one,
and target cardinality is uniformly at most one. Nevertheless a
computable bound on the necessary coefficient heights would contradict
Proposition~\ref{prop:singleton}, because the bounded-envelope theorem
would then produce the forbidden complete-list algorithm.
Indeed, writing a linear auxiliary as \(aX+bY+c\), its good
intersection with \(Y=0\) forces \(a\ne0\) and hence
\(t=-c/a\). If the reduced coefficient heights are at most \(H\),
then \(|a|\ge1/H\) and \(|c|\le H\), giving \(t\le H^2\).
Thus the missing coefficient bound would directly locate the halting
time. This example shows exactly what a degree-only cover omits.

\subsection{The Siegel and Faltings routes}

The inspected Siegel report separates a finite-tuple approximation
cover from a geometric bound on the number of curve points in each
covering subspace~\cite{repo-siegel}. Earlier framework packets replace
some abstract divisor arguments by explicit norm calculations.
The plane-elimination theorem broadens that geometric part: if a
subspace's annihilating function is represented by an auxiliary plane
polynomial \(G\), and \(\Good(F,G)\) holds, then a finite list in
that intersection has at most \(\deg F\deg G\) points. Neither
smoothness nor a special weight shape is needed. A rational function
presentation must first clear denominators and separately account for
points where those denominators vanish; that conversion is an
additional contract, not implicit in the theorem.

For an actually supplied finite global cover by such \(G\)'s, one can
now compute all the intersection lists and obtain a complete target
list after a decidable filter. More generally, a uniform degree and
coefficient envelope for pointwise auxiliaries suffices. This gives
a concrete alternative target for an arithmetic-geometric proof:
produce the finite envelope, rather than a distinguished cover.

What is not supplied is a proof that every curve in the general Siegel
or Faltings classes admits the required envelope with PR bounds.
The existing tuple-local covering proposal, even if proved, need not
imply that stronger assertion. Likewise, nonvanishing of a particular
resultant is a check on one exact auxiliary; it is not a construction
of an auxiliary for every target point. General positive height
bounds and complete rational-point algorithms for those classes
remain outside this report. No unprovability or independence claim
about those general arithmetic sentences follows from this limitation.

\subsection{Mordell--Weil descent}

The inspected Mordell--Weil packet already gives explicit contracting
height descent, coefficient extraction, and bounded selection from
\emph{supplied} quotient representatives~\cite{repo-mw}. We do not
present these again as new results. The reset-counter theorem can
organize a refinement process that alternates finitely many ranked
tasks, once its local PR transition and rank contracts are available.
The plane-intersection procedure can solve zero-dimensional rational
subproblems arising in a particular presentation. Neither construction
produces the missing arbitrary rational quotient representatives or
proves a uniform generation theorem. Those are separate mathematical
inputs, not a loop scheduling issue.

\subsection{Analytic names, Puiseux data, and exact coefficient tests}

The new arithmetic programs use exact rational coefficient arrays.
This restriction matters. The predicate \(\Good(F,G)\) compares
finitely many exact rational coefficients with zero. If coefficients
were supplied only by arbitrary Cauchy names, a negative result from
a finite precision test would not certify exact zero. For example,
\(F=Y^2-X^2-a\) and \(G=Y-X\) have elimination behavior depending
on whether \(a=0\). Treating an arbitrary name for \(a\) as an
exact rational code would import a zero test that the input does not
provide.

The Puiseux obstruction report makes the stronger point that exact
singular coefficient data can encode zero detection even when each
individual represented value happens to be rational~\cite{repo-puiseux}.
The current framework respects that obstruction. It supports exact
algebraic codes, approximate observations, or local constructions with
positive separation certificates. It does not infer exact singular
Puiseux coefficients from arbitrary Cauchy data, and it does not add
root multiplicity matching or a globally continuous branch label.

The demand compiler does resolve a different analytic difficulty:
iterating a fixed operation while its required input accuracy changes
can be performed with a first-order PR program. The correctness proof
fixes the final precision and uses the finite horizon invariant.
Whether the local analytic estimate is valid remains a mathematical
lemma to prove, just as before.

\subsection{Prime-distribution constants and all-future claims}

The prime-distribution packet gives exact arithmetic targets for
Siegel--Walfisz and Bombieri--Vinogradov with existential constants
\emph{before} the universal range of \(X\), and isolates the missing
effective uniform information~\cite{repo-primes}. The present
framework makes finite sums, rational approximation, and their
precision schedules easier to compile. It does not move a finite
search success at each tested \(X\) in front of the universal
quantifier. Nor does the reset theorem produce a stopping rank for
searching for a constant whose correctness is an unbounded
all-future property.

The same distinction applies to metastability. A program that returns
a stable finite window for each requested control function does not
thereby return a final convergence modulus. These are different output
contracts. Our new resource theorem proves termination only after
the finite rank condition itself has an internal proof; it is not a
general mechanism for converting an unbounded mathematical promise
into a computable certificate.

\subsection{A practical dependency table}

\begin{longtable}{>{\raggedright\arraybackslash}p{0.24\textwidth}>{\raggedright\arraybackslash}p{0.34\textwidth}>{\raggedright\arraybackslash}p{0.34\textwidth}}
\toprule
Difficulty & Available resolution here & Remaining input\\
\midrule\endhead
Changing precision inside recursion
& Explicit \code{PR 2} compiler; checked semantic laws and partial internal equations
& Finish the stated arithmetic coherence proofs in Lean.\\
Lower counters reset during descent
& Fixed-dimensional lexicographic fuel, including an automatic trajectory envelope
& Internal local invariant and strict-rank proofs for the application.\\
Curve fibre count lacks monic or weight form
& Bounded shear and Sylvester fibre multiplicity, giving the degree-product bound
& A verified nonzero resultant for the actual auxiliary equation.\\
Finite auxiliaries exist pointwise but no cover was selected
& Exhaust all auxiliaries within a supplied degree and coefficient envelope
& Proof of a uniform coefficient envelope.\\
Only a tuple-local cover or a cardinality bound is known
& Correct finite-list cardinality consequence
& Positive location data for a complete list or height bound.\\
Arbitrary Cauchy coefficient equality
& Approximate or apartness-guarded interfaces remain usable
& Exact input codes or a specific separation witness.\\
General MW, Siegel, Faltings, and prime estimates
& Explicit finite components and a narrower list of obligations
& The substantive general arithmetic or analytic theorem.\\
\bottomrule
\end{longtable}

\section{Lean definitions, checked results, and the completion path}
\label{sec:lean}

The archive contains a standalone Lean 4.32.0 development using \code{Std}
only. Ten modules were compiled, including the concrete definitions of
all main new programs and their unproved targets. Eleven selected theorem
declarations were audited with \texttt{\#print axioms}. Their dependencies
are contained in \code{propext} and \code{Quot.sound}; none uses
\code{sorryAx} or \code{Classical.choice}. These are facts about the
supplied Lean development, not a claim that all the report's application
theorems have already been formalized.

\subsection{The actual object-language backend}

The namespace \code{LeanProofs.ConciseIPRA} is the repository's namespace.
The reused standalone \code{ObjectKernel.lean} contains the actual
\code{Syntax}, \code{Deduction}, and \code{PR} declarations with their
imports arranged for \code{Std}. The bodies were compared with the
pinned current upstream sources after removing import lines and
boundary whitespace; the source audit records the hashes and equality
checks. The original MIT-0 license is included.

The central types and constructor interfaces are
\begin{lstlisting}
PR : Nat -> Type
PTm : Nat -> Type
Fm PTm n
PRAAx : Theory PTm
Prf PRAAx Gamma phi
\end{lstlisting}
Here \(n\) is the number of free variables. \code{PR k} has constructors
\code{zero}, \code{succ}, \code{proj}, \code{comp}, and \code{prec}.
The last has input symbols of arities \(k\) and \(k+2\), and returns
a symbol of arity \(k+1\). The step's arguments are the recursion index,
the accumulator, and the \(k\) parameters.

The formulas in \code{PRAAx} include the defining equations, successor
separation and injectivity, and quantifier-free induction. The supplied
\code{PRA.qfInd} rule expects a proof that the formula is
\code{Fm.QF}, a base proof, and a step proof in the correctly shifted
context. Equality and substitution are handled by
\code{Prf.eqCongr}, \code{Prf.eqTrans}, \code{Prf.subst}, and
\code{Prf.trans}. The code therefore uses the real proof relation
instead of introducing a new predicate named ``provable'' and assuming
its soundness~\cite{repo-pr}.

\subsection{Status at declaration level}

\begin{longtable}{>{\raggedright\arraybackslash}p{0.27\textwidth}>{\raggedright\arraybackslash}p{0.36\textwidth}>{\raggedright\arraybackslash}p{0.28\textwidth}}
\toprule
Module & Supplied result & Status and logical level\\
\midrule\endhead
\code{ObjectKernel}
& Arithmetic syntax and intuitionistic deduction
& Reused backend; compiled.\\
\code{BackwardDemand}
& Horizon, replay, and finite-query lemmas
& Reused unchanged; checked host theorems.\\
\code{TermCompiler}
& \code{symbol_correct}, \code{transport_equation}
& New checked \code{Prf PRAAx} constructions.\\
\code{DemandCompiler}
& \code{compiled : PR 2}; \code{denote_compiled}; recurrence laws
& New checked syntax and host semantic proofs.\\
\code{ObjectDemand}
& Generic backward zero and successor equations
& New checked \code{Prf PRAAx} constructions.\\
\code{ObjectTargets}
& Internal monus, horizon, replay, and full recurrence statements
& Compiled definitions of unproved propositions.\\
\code{LexFuel}
& PR fuel syntax, including a variable input parameter; checked denotation
& Syntax/semantics checked. Chain and termination theorems are targets.\\
\code{RationalRoots}
& Complete-root candidate program with zero-polynomial failure flag
& Executable definitions; correctness is an unproved target.\\
\code{EliminationTargets}
& Shear, Sylvester determinant, intersection and auxiliary-envelope programs
& Executable definitions; fourteen unproved target propositions.\\
\code{AxiomAudit}
& Dependency reports for eleven selected proved declarations
& Checked audit, not an additional mathematical theorem.\\
\bottomrule
\end{longtable}

All unproved targets are \code{def ... : Prop := ...}, not \code{axiom}
declarations or theorems closed by \code{sorry}. Type checking a target
verifies the well-formedness of its statement and definitions; it does
not prove that statement. The new program definitions can consequently
be inspected and tested without confusing an unfinished theorem with
a trusted assumption.

\subsection{Completing the internal demand theorem}

The first new compiler has an entirely internal correctness conclusion:
\begin{lstlisting}
symbol_correct (t : PTm k) (sigma : Fin k -> PTm n) :
  Prf PRAAx Gamma
    (Fm.eq (PTm.app (symbol t) sigma) (Syn.subst sigma t))
\end{lstlisting}
This theorem is already proved. The two-pass compiler is also a
concrete syntax constructor, but its complete internal correctness
still requires the following named targets.

\begin{enumerate}
\item \code{MonusSelfTarget}. Derive the predecessor and truncated
subtraction equations from \code{PRA.precZero} and
\code{PRA.precSucc}. Use equation induction to prove
\(N\dotminus N=0\) and the elementary successor cancellation laws.
\item \code{BackwardCoherenceTarget}. Use the already checked
\code{object_backward_zero} and \code{object_backward_succ} to prove
\eqref{eq:back-coherent} by QF induction on the reverse index. The
base and step have no order side condition.
\item \code{HorizonFinalTarget}. Substitute the monus-self equation
and the backward zero equation. Define a PR Boolean comparison symbol
and prove the guarded horizon-coherence equation for \(t\le N\).
\item \code{ForwardBaseTarget} and \code{ForwardStepTarget}. Unfold
the exact \code{PR.prec} symbols. Apply the term compiler to the
composite step body, just as in the completed backward proof.
\item Compare the two forward runs with the QF invariant
\(t\le N\to V(t;N+1,K)=V(t;N,D(N,K))\). This is the internal
replacement for the host theorem \code{advance_congr}.
\item \code{CompiledZeroTarget} and \code{CompiledSuccessorTarget}.
Unfold the final composition and apply the preceding equations.
The full human proof appears in Theorem~\ref{thm:demand}.
\end{enumerate}

Every step has a defined numerical term and an exact object formula.
There is no remaining appeal to an arbitrary host induction theorem or
an uninterpreted ``compiler correctness'' axiom. What remains is the
construction of the corresponding finite proof terms.

\subsection{Fixed-dimensional fuel: the uniformity distinction in types}

In \code{LeanProofs.ConciseIPRA.LexFuel},
\code{Rank d := Fin d -> Nat}; \code{lexDropB} implements the
lexicographic comparison. \code{run}, \code{coordMax},
\code{prefixMax}, and \code{trajectoryEnvelope} implement the finite
simulation and the envelope used in the human proof. The host
\code{fuel} definition is exactly \(F_d\).

The checked syntax constructors have types
\begin{lstlisting}
fuelSymbol      : Nat -> PR 1 -> PR 1
fuelParamSymbol : Nat -> PR 2 -> PR 2
boundSymbol     : Nat -> PR 2 -> PR 1
\end{lstlisting}
In the parameterized version, the input symbol \(h\) denotes
\(h(x,t)\), and the output denotes \(F_d(x,s)\). The checked theorem
\code{denote_boundSymbol} proves
\[
 \operatorname{evalB}(\code{boundSymbol}\ d\ h,x)
       =F_d(x,0).
\]
The input \(x\) remains a variable of one PR symbol for every fixed
\(d\); it has not been specialized separately to each numeral.
On the other hand, these constructors do not assert that evaluation of
their generated syntax is a PR function jointly in arbitrary \(d,x,s\).
The non-PR diagonal proved in Section~\ref{sec:descent} explains why.

The exact remaining targets are \code{FuelGrowthTarget},
\code{ControlledChainBoundTarget}, \code{BoundedTerminationTarget},
and \code{WitnessTarget}. Their proof order is:

\begin{enumerate}
\item Prove iterator monotonicity in its starting point and iteration
count. Then prove the fuel growth and monotonicity properties.
\item Implement bounded block boundaries and slices of a finite
controlled chain. The first coordinate drops at each boundary;
within a block apply the preceding-dimensional suffix theorem.
Accumulate the endpoint inequality by induction on the finite number
of blocks.
\item Prove initialization and preservation of the given Boolean
invariant along \code{run}, and persistence of inactivity for the
absorbing transition. An active state at time \(B\) yields the
forbidden \(B+1\)-state chain.
\item Apply the supplied exit implication to the resulting inactive
state. For the internal version, encode tuple and list data by
natural numbers and Booleanize every finite condition. Supply actual
object proofs of the particular transition's local assumptions.
\end{enumerate}

The host \code{Machine} structure deliberately contains data fields,
not unproved correctness axioms. Its arbitrary host functions must
still be replaced by fixed PR syntax before the internal theorem is
instantiated. In particular, well-foundedness of an abstract Lean
relation is not being used as an arithmetic termination certificate.

\subsection{Exact finite polynomial data and executable definitions}

The algebraic specification layer uses
\begin{lstlisting}
QPoly     := List Rat
PlanePoly := List (List Rat)
QPoint    := Rat x Rat
\end{lstlisting}
Coefficients increase with degree; row \(i\), column \(j\) is the
coefficient of \(X^iY^j\). Trailing zeros are permitted. Default reads
return zero. \code{eval} and \code{evalPlane} are Horner evaluation.
The arithmetic layer will encode these finite lists and canonical
rationals by natural numbers; the report's PR constructions provide
the required fixed programs.

\code{RationalRoots.lean} includes the actual denominator product,
integer coefficients, first/last nonzero coefficient selection,
bounded numerator and denominator boxes, evaluation filter, and
duplicate removal. Zero is included as a candidate and tested, which
handles initial powers of the variable. The zero polynomial returns
\code{none}; a nonzero polynomial with no rational roots returns
\code{some []}. Its unproved targets are
\code{RationalRootsTarget}, \code{DegreeBoundTarget}, and
\code{HeightBoundTarget}.

\code{EliminationTargets.lean} supplies actual programs for
\code{shear}, \code{goodShear}, \code{canonicalShear},
\code{sylvester}, \code{determinant}, \code{elimination},
\code{good}, \code{intersect}, \code{auxiliaryBox}, and
\code{envelopeList}. The determinant enumerates the \(N^N\) vectors
of length \(N\) over \(\{0,\ldots,N-1\}\), retains permutations,
computes their signs, and sums their polynomial products. This
inefficient but transparent definition makes its finite bound explicit.
The auxiliary box similarly enumerates finite coefficient vectors.
No polynomial operation in the target module is an opaque oracle.

The mathematical specifications are literal propositions, for example
\begin{lstlisting}
Good F G := good F G = true
common F G p := eval2 F p.1 p.2 = 0 /\ eval2 G p.1 p.2 = 0
\end{lstlisting}
and
\begin{lstlisting}
IntersectCompleteTarget :=
  forall F G, Good F G -> exists points,
    intersect F G = some points /\ points.Nodup /\
    (forall p, p in points <-> common F G p)
\end{lstlisting}
The displayed snippets use ASCII mathematical notation for readability;
the accompanying Lean source contains the checked syntax. The target is
not supplied as a theorem.

\subsection{Dependency order for the plane theorems}

A complete formalization should follow the mathematical dependencies,
which allow the complete-list theorem to be finished before the sharp
count theorem.

\begin{enumerate}
\item \textbf{Rational arithmetic and polynomial identities.}
Prove denominator clearing preserves evaluation, gcd cancellation,
the finite rational-root bound, and correctness of the root scan.
Develop coefficientwise addition, multiplication, evaluation, degree
support, and synthetic division for the finite arrays.
\item \textbf{Shear correctness.}
Prove the binomial coefficient identity for \code{shearCoeff}, then
\code{ShearEvalTarget}. A nonzero polynomial of degree at most
\(d+e\) cannot vanish at all \(d+e+1\) candidate directions; this
proves \code{CanonicalShearTarget} and constant leading coefficients.
\item \textbf{Elimination and complete lists.}
Prove determinant specialization and the adjugate identity. Apply
the evaluation vector \((1,v,\ldots,v^{N-1})\) to establish
\code{CommonGivesRootTarget}. Every specialized first equation has
the same nonzero leading coefficient, so its root routine succeeds.
These facts and the two finite membership scans prove
\code{IntersectCompleteTarget}, followed by
\code{IntersectHeightTarget}.
\item \textbf{Sharp cardinality.}
Prove \code{ResultantDegreeTarget} using the permutation weight
identity from Lemma~\ref{lem:elim-degree}. Prove the Vandermonde
determinant and column factor extraction to obtain
\code{FibreDivisibilityTarget}. Combine distinct linear-factor
powers and apply the degree bound, giving
\code{IntersectCardinalityTarget}.
\item \textbf{Auxiliary exhaustion.}
Prove \code{MonomialCountTarget}, \code{AuxiliaryBoxCountTarget},
and \code{AuxiliaryBoxCoverageTarget}. The last permits a padded
coefficient array different from the input array but coefficientwise
equal to it. Establish invariance of the polynomial programs under
that equality. Finite union membership and the supplied coverage
premise then give \code{EnvelopeCompleteTarget}; the length and
maximum lemmas give \code{EnvelopeCardinalityTarget} and
\code{EnvelopeHeightTarget}.
\end{enumerate}

For the embedded arithmetic layer, every polynomial identity becomes
a bounded conjunction over its finite coefficient support. The support
bound is produced by the same array operations. Every divisibility
argument carries its explicit quotient array. Rational cancellation
carries gcd or B\'ezout data when needed. In the enumeration proofs,
fix the arbitrary point code as a parameter and induct on processed
prefixes with a bounded membership invariant. A universal coverage
assumption remains in the fixed context and is instantiated after the
finite program invariant has been proved; it is never an induction
formula.

\subsection{The final native transfer and integration boundary}

The actual upstream theorem is
\begin{lstlisting}
pra_iff_native :
  (Prf PRAAx Gamma phi) <->
  NPrf (Gamma.map NFm.ofPRA) (NFm.ofPRA phi)
\end{lstlisting}
in \code{ConciseIPRA.NativeBiInterpretation}~\cite{repo-native}.
It is applied after the object proof has been constructed. It does
not accept an ambient equality of denotations in place of its left-hand
proof object.

To integrate, replace the standalone kernel import with
\code{import ConciseIPRA.PR}, retain the new modules in their own
namespaces, complete the pending object targets, and use this bridge
for the final theorem. The standalone development was compiled.
A build against the full Mathlib-backed repository dependency closure
was not run, and no remote repository files were modified.

\section{Reproduction, evidence, and remaining work}
\label{sec:validation}

\subsection{What the executable checks establish}

The Python companions use exact integers and \code{fractions.Fraction}
with the standard library. They are reference implementations and
finite regression diagnostics, not an iPRA proof checker. Their results
supplement the universal arguments in this report.

The lexicographic program checks 55 explicitly constructed maximal
controlled chains, 23 monotonicity comparisons, and eight resetting
counter runs. For \(h(t)=t+1\), the exact three-dimensional chain
starting at zero has 2,046 terms. The reset example beginning at
\((3,0,0)\) halts after 136 steps; its automatically derived
trajectory-envelope fuel is 267. Larger hierarchy values are
deliberately not evaluated; their mathematical treatment is symbolic.

The plane program verifies five exact classifications, including
nontrivial rational denominators and zero roots. It also performs
independent bounded rational-box scans and checks the unsupported
branches. The example
\[
             X^2+Y^2=5,\qquad XY=2
\]
returns exactly \((1,2),(2,1),(-1,-2),(-2,-1)\), while its canonical
elimination polynomial is \((U^2-9)^2\). The repeated fibres exercise
the part of the proof that a simple distinct-projection count misses.

The auxiliary-envelope example enumerates all 24 positive-degree
polynomials in the degree-one, coefficient-height-one box, before
discarding no good candidates. For the target
\(X^2+Y^2=1\) with the filter \(XY=0\), it returns the four axis
points and height bound one. Coverage in this example is justified
mathematically by the two auxiliaries \(X\) and \(Y\), both of which
are in the box; the execution does not claim to discover a proof of
coverage for an arbitrary target set.

\subsection{Rebuild commands}

From the archive's top directory:
\begin{lstlisting}
python3 verify.py
LEAN_BIN=/absolute/path/to/lean-4.32.0/bin/lean bash lean/build.sh
bash paper/build.sh
\end{lstlisting}
The Lean distribution and a normal TeX Live installation must be
available locally. The Python checks have no third-party dependency.
The supplied PDF can be read without installing either toolchain.
Build logs, the exact Lean version, per-claim status, source hashes,
and SHA-256 hashes of delivered files are included in the archive.

\subsection{What remains before the main results are kernel checked in iPRA}

The full object demand theorem still needs the internal arithmetic
coherence lemmas listed in Section~\ref{sec:lean}. The lexicographic
chain and bounded-termination theorems still need their Lean proofs,
followed by object-theory lowering for each fixed dimension. The
elimination, rational-root completeness, and auxiliary-envelope
theorems have fully specified concrete algorithms and complete human
arguments, but their general Lean correctness proofs and embedded
\code{Prf PRAAx} derivations are not supplied here.

These remaining tasks are different from the unresolved general
mathematical inputs. Proving the supplied algorithm's correctness in
Lean will not prove a uniform auxiliary coefficient bound for every
curve, locate arbitrary Mordell--Weil representatives, or produce the
missing effective analytic constants. Conversely, those open inputs
do not invalidate the finite theorems proved here. The framework
provides explicit interfaces at which a future result can be inserted
and its exact constructive consequence read off.

\begin{thebibliography}{99}
\bibitem{repo-root}
ProveIt repository. \emph{Primitive recursive arithmetic: root README.}
Snapshot \code{b0f6f535ffc16324683fef2c263260a88ae03755}.
\href{https://github.com/VladimirReshetnikov/ProveIt/blob/b0f6f535ffc16324683fef2c263260a88ae03755/Logic/PrimitiveRecursiveArithmetic/README.md}{Pinned source}.

\bibitem{repo-index}
ProveIt repository. \emph{iPRA proof-packet index}, including its verification caveat and 60-packet catalog.
\href{https://github.com/VladimirReshetnikov/ProveIt/blob/b0f6f535ffc16324683fef2c263260a88ae03755/Logic/PrimitiveRecursiveArithmetic/docs/proofs/README.md}{Pinned source}.

\bibitem{repo-pr}
ProveIt repository. \emph{ConciseIPRA arithmetic syntax and axioms},
\code{PR.lean}, with \code{Syntax.lean} and \code{Deduction.lean}.
\href{https://github.com/VladimirReshetnikov/ProveIt/blob/b0f6f535ffc16324683fef2c263260a88ae03755/Logic/PrimitiveRecursiveArithmetic/ConciseIPRA/Lean/ConciseIPRA/PR.lean}{Pinned source}.

\bibitem{repo-native}
ProveIt repository. \emph{Native bi-interpretation},
\code{NativeBiInterpretation.lean}; theorem \code{pra_iff_native}.
\href{https://github.com/VladimirReshetnikov/ProveIt/blob/b0f6f535ffc16324683fef2c263260a88ae03755/Logic/PrimitiveRecursiveArithmetic/ConciseIPRA/Lean/ConciseIPRA/NativeBiInterpretation.lean}{Pinned source}.

\bibitem{repo-demand}
ProveIt research packet. \emph{Mathematics in iPRA: finite demands, root supports, and norm certificates}.
\code{ipra_finite_demand}, especially Section 2 and
\code{lean/BackwardDemand.lean}.
\href{https://github.com/VladimirReshetnikov/ProveIt/tree/b0f6f535ffc16324683fef2c263260a88ae03755/Logic/PrimitiveRecursiveArithmetic/docs/proofs/foundations/ipra_finite_demand}{Pinned packet}.

\bibitem{repo-coherent}
ProveIt research packet. \emph{Coherent finite methods for intuitionistic primitive recursive arithmetic}.
\code{ipra_coherent_methods_2026-11-15}, especially the weighted norm theorem and standalone object kernel.
\href{https://github.com/VladimirReshetnikov/ProveIt/tree/b0f6f535ffc16324683fef2c263260a88ae03755/Logic/PrimitiveRecursiveArithmetic/docs/proofs/foundations/ipra_coherent_methods_2026-11-15}{Pinned packet}.

\bibitem{repo-observation}
ProveIt research packet. \emph{Finite observations and uniform constructions in iPRA}.
\code{ipra_finite_observation}.
\href{https://github.com/VladimirReshetnikov/ProveIt/tree/b0f6f535ffc16324683fef2c263260a88ae03755/Logic/PrimitiveRecursiveArithmetic/docs/proofs/foundations/ipra_finite_observation}{Pinned packet}.

\bibitem{repo-framework}
ProveIt research packet. \emph{Finite contracts for mathematics in intuitionistic primitive recursive arithmetic}.
\code{ipra_framework}.
\href{https://github.com/VladimirReshetnikov/ProveIt/tree/b0f6f535ffc16324683fef2c263260a88ae03755/Logic/PrimitiveRecursiveArithmetic/docs/proofs/foundations/ipra_framework}{Pinned packet}.

\bibitem{repo-embedding}
ProveIt research packet. \emph{Mathematics in iPRA through finite observations}.
\code{ipra_embedding_framework}, including the hyperelliptic and finite quotient-algebra arguments.
\href{https://github.com/VladimirReshetnikov/ProveIt/tree/b0f6f535ffc16324683fef2c263260a88ae03755/Logic/PrimitiveRecursiveArithmetic/docs/proofs/foundations/ipra_embedding_framework}{Pinned packet}.

\bibitem{repo-other}
ProveIt research packets. \emph{Certified Finitary Presentations},
\emph{Evidence-indexed finite mathematics},
\emph{iPRA certificate-and-modulus framework}, and
\emph{Uniform certificates for iPRA}.
Directories \code{ipra_embedding_framework_2026-11-15},
\code{ipra_evidence_framework}, \code{ipra_framework_2026-11-15},
and \code{ipra_uniform_certificates}, respectively.
\href{https://github.com/VladimirReshetnikov/ProveIt/tree/b0f6f535ffc16324683fef2c263260a88ae03755/Logic/PrimitiveRecursiveArithmetic/docs/proofs/foundations}{Pinned foundations directory}.

\bibitem{repo-geometry}
ProveIt research packet. \emph{Diophantine geometry in iPRA},
\code{diophantine_geometry_ipra/sections/02_heights.tex}, including the finite rational-root algorithm.
\href{https://github.com/VladimirReshetnikov/ProveIt/blob/b0f6f535ffc16324683fef2c263260a88ae03755/Logic/PrimitiveRecursiveArithmetic/docs/proofs/diophantine-number-theory/diophantine_geometry_ipra/sections/02_heights.tex}{Pinned source}.

\bibitem{repo-siegel}
ProveIt research packet. \emph{Siegel in iPRA: the general route and remaining obligations},
\code{siegel_ipra/paper/05_general_route.tex}.
\href{https://github.com/VladimirReshetnikov/ProveIt/blob/b0f6f535ffc16324683fef2c263260a88ae03755/Logic/PrimitiveRecursiveArithmetic/docs/proofs/diophantine-number-theory/siegel_ipra/paper/05_general_route.tex}{Pinned source}.

\bibitem{repo-faltings}
ProveIt research packet. \emph{Faltings' theorem and intuitionistic PRA},
\code{faltings_ipra/README.md}, with its positive height-bound and complete-list distinction.
\href{https://github.com/VladimirReshetnikov/ProveIt/blob/b0f6f535ffc16324683fef2c263260a88ae03755/Logic/PrimitiveRecursiveArithmetic/docs/proofs/diophantine-number-theory/faltings_ipra/README.md}{Pinned source}.

\bibitem{repo-mw}
ProveIt research packet. \emph{Mordell--Weil in iPRA},
\code{mordell_weil_ipra/paper/02_descent.tex} and
\code{03_elliptic.tex}.
\href{https://github.com/VladimirReshetnikov/ProveIt/tree/b0f6f535ffc16324683fef2c263260a88ae03755/Logic/PrimitiveRecursiveArithmetic/docs/proofs/diophantine-number-theory/mordell_weil_ipra}{Pinned packet}.

\bibitem{repo-primes}
ProveIt research packet. \emph{Prime distribution in iPRA},
\code{prime_distribution_ipra/paper/sections/07_sw_bv.tex}.
\href{https://github.com/VladimirReshetnikov/ProveIt/blob/b0f6f535ffc16324683fef2c263260a88ae03755/Logic/PrimitiveRecursiveArithmetic/docs/proofs/analytic-number-theory/prime_distribution_ipra/paper/sections/07_sw_bv.tex}{Pinned source}.

\bibitem{repo-puiseux}
ProveIt research packet. \emph{A uniform computability obstruction for arbitrary Cauchy-presented coefficients},
\code{puiseux_ipra/evidence/scope_obstruction.md}.
\href{https://github.com/VladimirReshetnikov/ProveIt/blob/b0f6f535ffc16324683fef2c263260a88ae03755/Logic/PrimitiveRecursiveArithmetic/docs/proofs/algebra/puiseux_ipra/evidence/scope_obstruction.md}{Pinned source}.

\bibitem{schmitz2014}
Sylvain Schmitz. \emph{Complexity Bounds for Ordinal-Based Termination}.
Reachability Problems 2014, LNCS 8762, pp.~1--19.
\href{https://arxiv.org/abs/1407.5896}{arXiv:1407.5896}.
Background on length bounds extracted from ordinal and lexicographic
termination arguments. Section~\ref{sec:descent} gives its own elementary
proof of the particular finite-dimensional bound used here.

\bibitem{sturmfels}
Bernd Sturmfels. \emph{Solving Systems of Polynomial Equations}.
CBMS Regional Conference Series in Mathematics, vol.~97, AMS, 2002.
Chapter 4, Resultants.
\href{https://math.berkeley.edu/~bernd/cbms.pdf}{Author-hosted text}.
Background for the Sylvester determinant; the degree, fibre, and
enumeration arguments used here are written out in full.
\end{thebibliography}

\end{document}
